\documentclass[10pt]{article}

\usepackage[utf8]{inputenc}
\usepackage[english]{babel}
\usepackage[margin=0.85in]{geometry}
\usepackage{amsmath,amsthm,amssymb,stmaryrd}
\usepackage{hyperref,thmtools,enumitem,xcolor}

\hypersetup{colorlinks,linkcolor={red!50!black},citecolor={blue!50!black},urlcolor={blue!80!black}}

\declaretheorem[name=Theorem,numberwithin=section]{theorem}
\declaretheorem[name=Lemma,sibling=theorem]{lemma}
\declaretheorem[name=Proposition,sibling=theorem]{proposition}
\declaretheorem[name=Corollary,sibling=theorem]{corollary}
\declaretheorem[name=Definition,sibling=theorem,style=definition]{definition}
\declaretheorem[name=Remark,sibling=theorem,style=remark]{remark}
\declaretheorem[name=Claim,sibling=theorem,style=remark]{claim}
\declaretheorem[name=Example,sibling=theorem,style=remark]{example}
\setlist[itemize]{parsep=0pt,topsep=2pt,before=\leavevmode}
\setlist[enumerate]{parsep=0pt,topsep=2pt,before=\leavevmode}
\setlength{\parindent}{1em}
\setlength{\parskip}{.15em}

\newcommand{\N}{\mathbb{N}}
\newcommand{\Prop}{\mathsf{Prop}}
\newcommand{\Type}{\mathsf{Type}}
\newcommand{\CIC}{\mathrm{CIC}}
\newcommand{\ZF}{\mathrm{ZF}}
\newcommand{\ZFC}{\mathrm{ZFC}}
\newcommand{\IZF}{\mathrm{IZF}}
\newcommand{\ZFh}{\mathrm{ZF}^{1/2}}
\newcommand{\Z}{\mathrm{Z}}
\newcommand{\KPP}{\mathrm{KP}(\mathcal P)}
\newcommand{\EM}{\mathrm{EM}}
\newcommand{\Con}{\mathrm{Con}}
\newcommand{\Ord}{\mathrm{Ord}}
\newcommand{\rk}{\operatorname{rank}}
\newcommand{\TC}{\operatorname{TC}}
\newcommand{\dom}{\operatorname{dom}}
\newcommand{\PP}{\mathcal{P}}
\newcommand{\Sat}{\mathrm{Sat}}
\newcommand{\WF}{\mathrm{WF}}
\newcommand{\rec}{\mathsf{rec}}
\newcommand{\wfrec}{\mathsf{wfrec}}
\newcommand{\uchoice}{\mathsf{uchoice}}
\newcommand{\sem}[1]{\llbracket #1\rrbracket}
\newcommand{\Lang}{\mathcal L}
\newcommand{\Terms}{\mathrm{Tm}}
\newcommand{\hartogs}{\aleph}
\newcommand{\M}{\mathfrak M}
\newcommand{\Nat}{\mathrm{Nat}}
\newcommand{\vl}{\mathrm{vl}}
\newcommand{\El}{\mathrm{El}}
\newcommand{\code}[1]{\ulcorner #1\urcorner}
\newcommand{\U}{\mathcal{U}}

\begin{document}

\title{$\ZFh$: Zermelo set theory with recursion,\\
  \large and why $\ZF$ proves its consistency}
\author{Mario Carneiro}
% Drafted with Claude; adjust authorship/acknowledgement as you see fit.
\date{Draft, \today}
\maketitle

\begin{abstract}
We define $\ZFh$, Zermelo set theory extended by a term language closed
under comprehension, images, and recursion along $\in$ and along
well-founded relations, together with the axiom that every term has an
image. This is the set theory that the sets-as-trees interpretation of the
Calculus of Inductive Constructions with two universes and excluded middle
satisfies: it has full Separation and Powerset, well-founded recursion for
every term-definable operation, and Replacement only for term-definable
functions, never for functional relations. We show that $\ZF$ proves the
consistency of $\ZFh$. The model is $V_\lambda$ for an ordinal $\lambda$
of cofinality $\omega$, closed under a countable family of definable
ordinal functions that majorise the ranks of all terms. The construction
extends Miquel's observation that Zermelo set theory already proves
Replacement for its Skolem terms because each term carries a syntactic
rank bound; recursion turns finite bounds into transfinite ones, and the
absence of a description operator is what keeps the bounds computable
from the shape of a term. We also give a term-free axiomatisation, show
that $V_\lambda$ refutes Replacement, and place $\ZFh$ strictly between
$\Z$ and $\ZF$. We then prove that $\CIC_2+\EM$ interprets $\ZFh$, the
only non-standard ingredient being that well-founded recursion along a
set relation is obtained from $\mathsf{Acc}$-recursion without choosing
representatives, and that the type theory proves $\Con(\ZFh)$ outright.
Adding universes on both sides, $\CIC_\omega+\EM$ interprets $\ZFh_\omega$,
$\ZFh$ with an $\omega$-tower of universes, and $\ZF$ proves
$\Con(\ZFh_\omega)$ by the same majorant argument, with all universes
interpreted as $V_{\lambda_i}$ for closed ordinals $\lambda_i$.
\end{abstract}

\section{Introduction}

Type theories with an impredicative $\Prop$ and a hierarchy of universes,
such as Coq's or Lean's, interpret set theory by sets-as-trees: a set is a
well-founded tree whose branching types live in a universe
\cite{aczel78,werner97,barras10}. The interpretation validates Zermelo's
axioms with full Separation, and it validates recursion along $\in$ and
along any well-founded relation. What it does not validate is Replacement:
from $\forall x{\in}a\,\exists!y\,\varphi(x,y)$ one cannot form the image,
because the $\exists$ is proof-irrelevant and yields no function. Werner
and Barras obtain Replacement only by adding a description or unique
choice axiom \cite{werner97,barras10}, and the assembly models of
\cite{vdbergotten23} refute such axioms. So the set theory of the type
theory has Replacement for \emph{functions the theory can write down}, and
nothing more.

This note isolates that set theory, which we call $\ZFh$, and proves that
$\ZF$ proves its consistency. The idea comes from Miquel
\cite{miquel06}, who observed that in a Skolemised presentation of
Zermelo set theory, with terms for pairing, union, powerset, $\omega$, and
comprehension, the image $\{t(x):x\in u\}$ of any term $t$ is already
definable, because a bounding set $B(t(x),x\in u)$ can be computed by
structural recursion on $t$; he used this to show that his four-sorted
pure type system $\lambda\mathrm Z$ is equiconsistent with Zermelo set
theory. Comprehension is rank-neutral, so the bound never looks at the
formulas inside a term, only at its shape. Once recursion is added, the
bounds are no longer terms but ordinal functions, still computed from the
shape of the term. Take $\lambda$ closed under all of them; then
$V_\lambda$ satisfies every axiom of $\ZFh$, and $\ZF$ proves that such a
$\lambda$ exists by a countable iteration, with no regularity or closure
under powerset required of $\lambda$ beyond being a limit.

\paragraph{Main results.}
\begin{itemize}
\item $\ZF\vdash\Con(\ZFh)$ (Theorem~\ref{thm:main}); hence $\ZFh\nvdash\Con(\ZF)$
  and $\ZFh$ is a proper subtheory of $\ZF$, since $V_\lambda$ refutes
  Replacement (Proposition~\ref{prop:norepl}).
\item $\ZFh$ has a term-free axiomatisation: $\Z$ plus $\in$-induction,
  ``every set has a rank and every $V_\alpha$ exists,'' and the totality
  of a countable family of definable ordinal functions
  (Proposition~\ref{prop:termfree}).
\item The recursion-free fragment of $\ZFh$ is conservative over $\Z$
  \cite{miquel06}; $\ZFh$ itself proves $V_{\omega+\omega}$ exists, so
  recursion is exactly what separates it from $\Z$.
\end{itemize}
\begin{itemize}
\item $\CIC_2+\EM$ interprets $\ZFh$ by sets-as-trees
  (Theorem~\ref{thm:interp}), and proves $\Con(\ZFh)$, so the
  interpretation cannot be reversed (Proposition~\ref{prop:con}).
\item With universes: $\CIC_{n+2}+\EM$ interprets $\ZFh_n$ and
  $\CIC_\omega+\EM$ interprets $\ZFh_\omega$ (Theorem~\ref{thm:interpn}),
  and $\ZF\vdash\Con(\ZFh_\omega)$ (Theorem~\ref{thm:univ}).
\end{itemize}
What is missing for $\ZF\vdash\Con(\CIC_\omega+\EM)$ is the converse
model, of $\CIC_{n+2}+\EM$ inside $\ZFh_{n+1}$, and
Section~\ref{sec:univ} explains why the construction of this note is the
template for it: the type theory's function spaces must be interpreted
by term-definable functions, which carry rank majorants, rather than by
all set functions, which do not.

\section{The theory}\label{sec:theory}

\subsection{Language}

$\ZFh$ is a first-order theory with equality in a language $\Lang$ whose
terms and formulas are defined simultaneously. Variables, $\omega$,
$\{t,u\}$, $\bigcup t$, $\PP(t)$ are terms; $t=u$ and $t\in u$ are atomic
formulas; formulas are closed under $\wedge,\vee,\to,\bot$ and the
unbounded quantifiers $\forall x,\exists x$. In addition there are four
variable-binding term formers:
\begin{align*}
&\{x\in t\mid\varphi\} && \text{comprehension ($x$ bound in $\varphi$, not in $t$)}\\
&\{t\mid x\in u\} && \text{image ($x$ bound in $t$, not in $u$)}\\
&\rec_{x,g.\,G}(t) && \text{$\in$-recursion ($x,g$ bound in $G$)}\\
&\wfrec_{x,g.\,G}(X,R,t) && \text{recursion along a relation ($x,g$ bound in $G$)}
\end{align*}
Every term former may carry further free variables as parameters. Terms
are finite syntax trees, so $\Terms$, the set of terms, is countable and
is a set in any metatheory containing arithmetic. Writing $\vec p$ for the
parameters of a term, the intended reading is: $\{t\mid x\in u\}$ is the
image of $u$ under $x\mapsto t$; $\rec_{x,g.G}(a)$ is $r(a)$ for the
function $r$ on $\TC(\{a\})$ with $r(x)=G(x,r{\restriction}x)$, where
$r{\restriction}x=\{\langle y,r(y)\rangle:y\in x\}$; and
$\wfrec_{x,g.G}(X,R,a)$ is $r(a)$ for the function $r$ on $X$ with
$r(x)=G(x,\{\langle y,r(y)\rangle:y\in X,\ \langle y,x\rangle\in R\})$,
provided $R$ is well-founded on $X$.

Formally one may regard each binding former as a family of ordinary
function symbols, one for each choice of the bound syntax: $C_\varphi$ of
arity $|\vec p|+1$, $I_t$, $\mathrm{Rec}_G$, $\mathrm{Wf}_G$. Then $\Lang$ is
an ordinary countable first-order signature, defined inductively together
with its formulas, and $\ZFh$ is an ordinary first-order theory; this is
the presentation we use for the model.

\subsection{Axioms}

$\ZFh$ has classical first-order logic with equality and the following
axioms, where $\WF(X,R)$ abbreviates
$\forall S\,[S\subseteq X\wedge S\ne\emptyset\to\exists x{\in}S\,\forall y{\in}S\,\langle y,x\rangle\notin R]$,
$\Nat(n)$ is Miquel's formula ``$n$ belongs to every set containing
$\emptyset$ and closed under successor,'' and $\langle y,z\rangle$ is the
Kuratowski pair.
\begin{align*}
(\mathrm{Ext}) &\quad \forall a\,b\,[\forall x\,(x\in a\leftrightarrow x\in b)\to a=b]\\
(\mathrm{Pair}) &\quad \forall a_1\,a_2\,x\,[x\in\{a_1,a_2\}\leftrightarrow x=a_1\vee x=a_2]\\
(\mathrm{Union}) &\quad \forall a\,x\,[x\in\textstyle\bigcup a\leftrightarrow\exists y\,(y\in a\wedge x\in y)]\\
(\mathrm{Pow}) &\quad \forall a\,x\,[x\in\PP(a)\leftrightarrow x\subseteq a]\\
(\mathrm{Inf}) &\quad \forall x\,[x\in\omega\leftrightarrow\Nat(x)]\\
(\mathrm{Comp}) &\quad \forall\vec p\,a\,x\,[x\in\{z\in a\mid\varphi\}\leftrightarrow x\in a\wedge\varphi(x)]\\
(\mathrm{Img}) &\quad \forall\vec p\,u\,y\,[y\in\{t\mid x\in u\}\leftrightarrow\exists x\,(x\in u\wedge y=t)]\\
(\mathrm{Rec}) &\quad \forall\vec p\,a\,[\rec_{x,g.G}(a)=G(a,\{\langle y,\rec_{x,g.G}(y)\rangle\mid y\in a\})]\\
(\mathrm{WfRec}) &\quad \forall\vec p\,X\,R\,a\,[\WF(X,R)\wedge a\in X\to
  \wfrec_{x,g.G}(X,R,a)=G(a,\{\langle y,\wfrec_{x,g.G}(X,R,y)\rangle\mid y\in R^{-1}[a]\})]\\
(\in\text{-}\mathrm{Ind}) &\quad \forall\vec p\,[\forall x\,(\forall y{\in}x\,\varphi(y)\to\varphi(x))\to\forall x\,\varphi(x)]
\end{align*}
Here $G(a,s)$ denotes $G$ with $x,g$ replaced by $a,s$,
$R^{-1}[a]$ abbreviates the comprehension term $\{y\in X\mid\langle y,a\rangle\in R\}$,
and the inner images in (Rec) and (WfRec) are terms of $\Lang$. All
schemata range over all terms and formulas of $\Lang$, including those
that themselves contain recursion.

\begin{remark}
(i) There is no description operator $\iota y.\varphi$ and no choice
operator. Adding one for provably functional $\varphi$ would turn (Img)
into full Replacement, so the content of $\ZFh$ lies entirely in the
choice of term formers. Each former present has a rank bound computable
from its shape; description is the one natural former that has none.
(ii) The fragment without $\rec$ and $\wfrec$ is Miquel's $\Z^{\mathrm{sk}}$
plus (Img), and (Img) is a theorem there \cite[Prop.~6]{miquel06}, so that
fragment is conservative over $\Z$ \cite[Prop.~5]{miquel06}.
(iii) $\ZFh$ proves that $\TC(a)$ exists (by $\rec$), that every set has
a rank, that $V_\alpha$ exists for every ordinal $\alpha$ (by $\rec$ on
ordinals), and hence that $V_{\omega+\omega}$ exists; none of these is
provable in $\Z$ \cite{esserhinnion99,mathias01}.
(iv) (WfRec) subsumes recursion on $\omega$ and on any well-ordered set,
and (Rec) is the special case $R={\in}{\restriction}\TC(\{a\})$; we keep
both because they correspond to different constructs on the type-theoretic
side.
\end{remark}

\section{Rank majorants}\label{sec:rank}

We work in $\ZF$. Write $\rk(x)$ for the von Neumann rank and
$\hartogs(S)$ for the Hartogs number of $S$, the least ordinal that does
not inject into $S$. We use the elementary bounds
$\rk\{a,b\}=\max(\rk a,\rk b)+1$, $\rk\bigcup a\le\rk a$,
$\rk\PP(a)=\rk a+1$, $\rk\{z\in a\mid\ldots\}\le\rk a$,
$\rk\langle a,b\rangle\le\max(\rk a,\rk b)+2$, and, for a function $r$
with domain $d$, $\rk\{\langle y,r(y)\rangle:y\in d\}\le\max(\rk d,\sup_{y\in d}\rk r(y))+3$.

\begin{definition}[majorants]\label{def:maj}
To each term $t$ assign an ordinal function $f_t:\Ord\to\Ord$ by
recursion on $t$. The argument $\alpha$ is a bound on the ranks of the
values of all free variables of $t$.
\begin{align*}
f_x(\alpha)&=\alpha & f_\omega(\alpha)&=\omega\\
f_{\{t,u\}}(\alpha)&=\max(f_t(\alpha),f_u(\alpha))+1 & f_{\bigcup t}(\alpha)&=f_t(\alpha)\\
f_{\PP(t)}(\alpha)&=f_t(\alpha)+1 & f_{\{x\in t\mid\varphi\}}(\alpha)&=f_t(\alpha)\\
f_{\{t\mid x\in u\}}(\alpha)&=f_t(\max(\alpha,f_u(\alpha)))+1
\end{align*}
For $\rec_{x,g.G}(t)$, define $h:\Ord\to\Ord$ by transfinite recursion,
$h(\gamma)=f_G\bigl(\max(\alpha,\gamma,\sup_{\delta<\gamma}h(\delta)+3)\bigr)$,
and set $f_{\rec_{x,g.G}(t)}(\alpha)=\sup_{\gamma\le f_t(\alpha)}h(\gamma)$.
For $\wfrec_{x,g.G}(X,R,t)$, with $\beta=f_X(\alpha)$, define
$h(\gamma)=f_G\bigl(\max(\alpha,\beta,\sup_{\delta<\gamma}h(\delta)+3)\bigr)$
and set $f_{\wfrec_{x,g.G}(X,R,t)}(\alpha)=\sup_{\gamma<\hartogs(V_{\beta+1})}h(\gamma)$.
\end{definition}

The assignment $t\mapsto f_t$ is a single $\ZF$-definable function of
two arguments $(t,\alpha)$: it is a recursion on the syntax tree of $t$
whose clauses use ordinal arithmetic, suprema over sets, transfinite
recursion on ordinals, and Hartogs numbers, all of which are
$\ZF$-definable total operations. It never refers to the truth of any
formula, because comprehension is rank-neutral.

\begin{definition}[evaluation]\label{def:eval}
Fix a limit ordinal $\lambda>\omega$. By simultaneous recursion on syntax
define, for terms $t$ and assignments $\rho$ of sets to variables, a set
$\sem t_\rho$, and for formulas $\varphi$ a truth value $\Sat_\lambda(\varphi,\rho)$:
\begin{itemize}
\item $\sem x_\rho=\rho(x)$; $\sem\omega_\rho=\omega$; pairing, union, and
  powerset are interpreted by the set-theoretic operations in $V$.
\item $\sem{\{x\in t\mid\varphi\}}_\rho=\{z\in\sem t_\rho:\Sat_\lambda(\varphi,\rho[x{:=}z])\}$
  and $\sem{\{t\mid x\in u\}}_\rho=\{\sem t_{\rho[x:=z]}:z\in\sem u_\rho\}$.
\item $\sem{\rec_{x,g.G}(t)}_\rho=r(\sem t_\rho)$ for the unique
  function $r$ on $\TC(\{\sem t_\rho\})$ with
  $r(z)=\sem G_{\rho[x:=z,\,g:=r\restriction z]}$, which exists by the
  recursion theorem of $\ZF$; likewise for $\wfrec$ when
  $\WF(\sem X_\rho,\sem R_\rho)$ and $\sem t_\rho\in\sem X_\rho$, and
  $\emptyset$ otherwise.
\item $\Sat_\lambda(t\in u,\rho)$ iff $\sem t_\rho\in\sem u_\rho$;
  connectives as usual; $\Sat_\lambda(\forall x\,\varphi,\rho)$ iff
  $\Sat_\lambda(\varphi,\rho[x{:=}z])$ for all $z\in V_\lambda$.
\end{itemize}
\end{definition}
The evaluation is total on all assignments into $V$ and only the
quantifier clause mentions $\lambda$. For $\rho$ into $V_\lambda$ and terms
whose values stay in $V_\lambda$ it is the value of $t$ in the
$\Lang$-structure $\M_\lambda$ with universe $V_\lambda$ whose function
symbols are interpreted by these clauses.

\begin{lemma}[rank lemma]\label{lem:rank}
For every term $t$, every $\lambda$, and every assignment $\rho$ with
$\rk\rho(y)\le\alpha$ for all free $y$ of $t$:
$\rk\sem t_\rho\le f_t(\alpha)$.
\end{lemma}
\begin{proof}
Induction on $t$. The constant, pairing, union, powerset, and
comprehension cases are the elementary bounds; note that comprehension
gives $\rk\le\rk\sem t_\rho$ whatever $\Sat_\lambda$ says. For the image,
each $z\in\sem u_\rho$ has rank $<f_u(\alpha)$, so
$\rk\sem t_{\rho[x:=z]}\le f_t(\max(\alpha,f_u(\alpha)))$ by induction,
and a set of such elements has rank at most one more. For $\rec$, let
$a=\sem t_\rho$ and $r$ the recursion function. We show
$\rk r(z)\le h(\rk z)$ for $z\in\TC(\{a\})$ by $\in$-induction on $z$: the
assignment $\rho[x:=z,g:=r{\restriction}z]$ has all values of rank at most
$\max(\alpha,\rk z,\sup_{y\in z}\rk r(y)+3)\le\max(\alpha,\gamma,\sup_{\delta<\gamma}h(\delta)+3)$
where $\gamma=\rk z$, using $\rk y<\gamma$ and the induction hypothesis;
apply the induction hypothesis for $G$. Then
$\rk r(a)\le h(\rk a)\le\sup_{\gamma\le f_t(\alpha)}h(\gamma)$. For
$\wfrec$ the same induction runs along the $R$-rank $\varrho(z)$ of
$z\in\sem X_\rho$, with $\rk z<\beta$ for all $z$, giving
$\rk r(z)\le h(\varrho(z))$. The $R$-rank is a surjection of
$\sem X_\rho$ onto an ordinal $\theta$, so $\theta$ injects into
$\PP(\sem X_\rho)\subseteq V_{\beta+1}$ by taking fibres, whence
$\theta<\hartogs(V_{\beta+1})$ without any use of choice.
\end{proof}

\section{The model}\label{sec:model}

\begin{definition}
An ordinal $\lambda$ is \emph{$\Lang$-closed} if it is a limit ordinal
$>\omega$ and $f_t(\alpha)<\lambda$ for every term $t$ and every $\alpha<\lambda$.
\end{definition}

\begin{lemma}\label{lem:closed}
$\ZF$ proves that $\Lang$-closed ordinals exist, and indeed that they are
unbounded.
\end{lemma}
\begin{proof}
Given $\alpha_0$, put $\lambda_0=\max(\alpha_0,\omega)$ and
$\lambda_{n+1}=\sup\{f_t(\alpha):t\in\Terms,\ \alpha\le\lambda_n\}+1$.
Each $\lambda_{n+1}$ is a set-supremum, because $(t,\alpha)\mapsto f_t(\alpha)$
is a definable function on the set $\Terms\times(\lambda_n+1)$, so
Replacement applies. Since $f_x(\alpha)=\alpha$, the sequence is strictly
increasing, and $\lambda=\sup_n\lambda_n$ is a limit ordinal. If
$\alpha<\lambda$ then $\alpha\le\lambda_n$ for some $n$ and
$f_t(\alpha)<\lambda_{n+1}\le\lambda$.
\end{proof}

Note that $\lambda$ has cofinality $\omega$ and is not required to be a
cardinal, a strong limit, or regular.

\begin{theorem}\label{thm:main}
$\ZF$ proves: for every $\Lang$-closed $\lambda$, $\M_\lambda\models\ZFh$.
Consequently $\ZF\vdash\Con(\ZFh)$.
\end{theorem}
\begin{proof}
First, $\M_\lambda$ is a genuine $\Lang$-structure: for $\rho$ into
$V_\lambda$, every $\sem t_\rho$ lies in $V_\lambda$, by
Lemma~\ref{lem:rank} and closure, so the clauses of
Definition~\ref{def:eval} interpret each function symbol as an operation
on $V_\lambda$. We check the axioms; each check is uniform in the axiom
scheme, so the argument is a single $\ZF$ proof of
``$\forall\sigma\in\ZFh.\ \M_\lambda\models\sigma$,'' and the
consequence follows by the soundness of first-order logic for set-sized
structures, which $\ZF$ proves.

(Ext), (Pair), (Union), (Inf): $V_\lambda$ is transitive, closed under
pairs and unions, and contains $\omega$; inductive sets in $V_\lambda$
contain $\omega$, so $\Nat^{\M_\lambda}(x)$ iff $x\in\omega$.
(Pow): $\PP(a)\in V_\lambda$ for $a\in V_\lambda$ since $\lambda$ is a
limit, and every subset of $a$ lies in $V_\lambda$, so the powerset in
$V$ is the powerset in $\M_\lambda$.
(Comp) and (Img) hold by the evaluation clauses, given that the
right-hand sides quantify over $V_\lambda$ and that the elements
$z\in\sem u_\rho$ and values $\sem t_{\rho[x:=z]}$ lie in $V_\lambda$.
(Rec) and (WfRec) hold by the defining property of the recursion
function $r$, all of whose values lie in $V_\lambda$ by
Lemma~\ref{lem:rank}. The hypothesis $\WF(X,R)$ is absolute between
$\M_\lambda$ and $V$, because every subset of $\sem X_\rho$ belongs to
$V_\lambda$.
($\in$-Ind): in any transitive set, and for any assignment, the
relativised induction scheme holds because $\in$ is well-founded in the
metatheory: if $\Sat_\lambda(\varphi,\rho[x{:=}z])$ failed for some
$z\in V_\lambda$, an $\in$-minimal such $z$ would violate the hypothesis.
Classical logic holds in every structure.
\end{proof}

\begin{corollary}\label{cor:strict}
If $\ZF$ is consistent then $\ZFh\nvdash\Con(\ZF)$, and $\ZFh$ does not
interpret $\ZF$.
\end{corollary}
\begin{proof}
If $\ZFh\vdash\Con(\ZF)$ then, arguing in $\ZF$: $\neg\Con(\ZF)$ gives a
$\ZF$-proof of $\bot$, hence by $\Sigma^0_1$-completeness
$\ZFh\vdash\neg\Con(\ZF)$, so $\ZFh$ is inconsistent, contradicting
Theorem~\ref{thm:main}. Thus $\ZF\vdash\Con(\ZF)$, which is impossible.
An interpretation of $\ZF$ in $\ZFh$ would give $\ZFh\vdash\Con(\ZF)$
by formalising it.
\end{proof}

\begin{proposition}\label{prop:norepl}
$\M_\lambda\not\models$ Replacement, for $\lambda$ constructed as in
Lemma~\ref{lem:closed} from $\alpha_0=\omega$. Hence $\ZFh\nvdash$
Replacement and $\ZFh\subsetneq\ZF$.
\end{proposition}
\begin{proof}[Proof sketch]
The function $n\mapsto\lambda_n$ is definable in $\M_\lambda$: the set
$\Terms$ is hereditarily finite, and for $\alpha<\lambda$ the computation
of $f_t(\alpha)$ consists of transfinite recursions along ordinals below
$\lambda$ and Hartogs numbers of sets $V_{\beta+1}$ with $\beta<\lambda$,
whose values are below $\lambda$ by closure; the witnessing functions are
sets of pairs of ordinals below $\lambda$, hence elements of $V_\lambda$,
and the recursion clauses are absolute for transitive sets containing the
relevant powersets. So $\M_\lambda$ defines a functional relation on
$\omega$ whose image is cofinal in $\lambda$, and the image is not an
element of $V_\lambda$.
\end{proof}

\section{A term-free axiomatisation and placement}\label{sec:termfree}

Write $\mathrm{Rk}$ for the statement ``every set $a$ has a rank: there
is an ordinal $\alpha$ and a function $\beta\mapsto V_\beta$ on $\alpha+1$
satisfying $V_\beta=\bigcup_{\gamma<\beta}\PP(V_\gamma)$, with
$a\in V_\alpha$.'' For each term $t$ write $\mathrm{Tot}(f_t)$ for the
statement ``for every $\alpha$ the value $f_t(\alpha)$ exists,'' where
``$y=f_t(\alpha)$'' is the first-order formula asserting the existence of
witnessing functions for the recursions in Definition~\ref{def:maj}, with
Hartogs numbers defined in the usual way.

\begin{proposition}\label{prop:termfree}
$\ZFh$ is a conservative extension of the theory
$\ZFh_0:=\Z+{\in}\text{-}\mathrm{Ind}+\mathrm{Rk}+\{\mathrm{Tot}(f_t):t\in\Terms\}$
in the language of $\in$, and $\ZFh_0$ is interpretable in $\ZFh$.
\end{proposition}
\begin{proof}[Proof sketch]
$\ZFh$ proves $\mathrm{Rk}$ and each $\mathrm{Tot}(f_t)$ using $\rec$ and
$\wfrec$, since every operation in Definition~\ref{def:maj} is a
$\wfrec$ along an ordinal or a comprehension. Conversely, in $\ZFh_0$ every
term of $\Lang$ can be eliminated: interpret $t$ by the formula
``$y$ is the value of $t$,'' where for the recursion formers the value is
defined via ``there is an attempt,'' a function on a transitive subset of
$\TC(\{a\})\times V_{f(\ldots)}$ satisfying the clauses, with $f$ the
majorant. Attempts exist and are unique by $\in$-induction, using
Separation from $\PP(\TC(\{a\})\times V_\beta)$ to collect them, which is
where $\mathrm{Rk}$ and $\mathrm{Tot}$ supply the bound $\beta$; $\TC(\{a\})$
itself is the least transitive subset of $V_{\rk a+1}$ containing $a$.
With terms eliminated, (Img) is an instance of Separation from
$V_{f_t(\rk u)+1}$, and the remaining axioms are their own translations.
The retraction is the identity on $\in$-formulas, giving conservativity
as in \cite[\S2.3]{miquel06}.
\end{proof}

So $\ZFh$ is Zermelo set theory plus $\in$-induction, ranks, and the
totality of a countable family of definable ordinal functions closed
under successor, suprema, transfinite iteration, and Hartogs numbers of
the cumulative hierarchy. Its proof-theoretic content is the size of the
least ordinal closed under this family, an ordinal of cofinality
$\omega$.

\begin{remark}[placement]
Power Kripke--Platek set theory $\KPP$ \cite{rathjen14} proves the
$\Sigma^{\PP}$-recursion theorem and $\Sigma^{\PP}$-Replacement, so
$\Z+\KPP$, i.e.\ Zermelo with full Separation and
$\Delta_0^{\PP}$-Collection, proves every axiom of $\ZFh_0$; and
$V_\lambda\models\Z+\KPP$ for a club of $\lambda$ by L\'evy reflection.
Thus $\Z\subsetneq\ZFh\subseteq\Z+\KPP\subsetneq\ZF$, all provably in
$\ZF$. We do not expect $\ZFh$ to prove $\Delta_0^{\PP}$-Collection,
since Collection bounds witnesses that no term can locate. The
primitive recursive set functions of Jensen and Karp \cite{jensenkarp71},
with powerset added, are exactly the recursion-free-of-formulas part of
our term language, and their totality is the core of $\mathrm{Tot}$.
\end{remark}

\section{$\CIC_2+\EM$ interprets $\ZFh$}\label{sec:tt}

\subsection{The type theory}

$\CIC_2+\EM$ is the extensional theory of \cite[\S2]{paper1}: an
impredicative $\Prop$, predicative universes $\Prop:\Type_0:\Type_1$ with
$\Type_1$ the top sort, $\Pi$, $\Sigma$, $W$, finite types, $\N$,
$=$ with reflection, propositional truncation eliminating into $\Prop$,
quotients, $\mathsf{propext}$, $\mathsf{em}$, and in addition an
accessibility predicate. For $X:s$ and $R:X\to X\to\Prop$:
\[
\frac{}{\Gamma\vdash\mathsf{Acc}_R:X\to\Prop}\qquad
\frac{\Gamma\vdash h:\Pi y{:}X.\,R\,y\,x\to\mathsf{Acc}_R\,y}{\Gamma\vdash\mathsf{acc}\,x\,h:\mathsf{Acc}_R\,x}
\]
\[
\frac{\Gamma,x{:}X\vdash C:s'\quad
      \Gamma\vdash f:\Pi x{:}X.\,(\Pi y{:}X.\,R\,y\,x\to C[y/x])\to C}
     {\Gamma\vdash\mathsf{ind}^{\mathsf{Acc}}_C f:\Pi x{:}X.\,\mathsf{Acc}_R\,x\to C}
\]
with $\mathsf{ind}^{\mathsf{Acc}}_C f\,x\,(\mathsf{acc}\,x\,h)\equiv
f\,x\,(\lambda y\,r.\,\mathsf{ind}^{\mathsf{Acc}}_C f\,y\,(h\,y\,r))$, the
motive $C$ in any sort. This is Coq's $\mathsf{Acc}$ with its singleton
elimination and Lean's $\mathsf{WellFounded.fix}$; in the assembly models
it is realized by a fixed-point combinator whose termination on
realizers of $x$ is proved by induction on $\mathsf{Acc}_R\,x$.

\subsection{Sets as trees}

Following Aczel and Werner \cite{aczel78,werner97}, let
$\mathsf{set}:=Wx{:}\Type_0.\,x$, a type in $\Type_1$. For
$a=\mathsf{tree}\,X\,f$ write $\bar a:=X$ and $f_a:=f:\bar a\to\mathsf{set}$.
Bisimulation and membership are defined by $W$-recursion into
$\mathsf{set}\to\Prop$, which is a type in $\Type_1$:
\[
a\approx b:=(\Pi i{:}\bar a.\,\exists j{:}\bar b.\,f_a\,i\approx f_b\,j)\wedge
            (\Pi j{:}\bar b.\,\exists i{:}\bar a.\,f_a\,i\approx f_b\,j),\qquad
a\in b:=\exists j{:}\bar b.\,a\approx f_b\,j.
\]
$\approx$ is an equivalence relation and $\in$ respects it; extensional
equality of sets is $\approx$, so the interpretation of $\ZFh$ below
reads $=$ as $\approx$. (With quotient types one may instead pass to
$\mathsf{set}/{\approx}$; nothing changes.) Kuratowski pairs, $\subseteq$,
and $\WF$ are defined from $\in$ in the usual way.

\begin{definition}[interpretation of $\Lang$]\label{def:star}
To each term $t$ of $\Lang$ with free variables $\vec p$ assign a term
$t^*:\mathsf{set}^{|\vec p|}\to\mathsf{set}$, and to each formula
$\varphi$ a term $\varphi^*:\mathsf{set}^{|\vec p|}\to\Prop$, by recursion
on syntax:
\begin{align*}
\omega^* &:= \mathsf{tree}\,\N\,(n\mapsto\underline n)\ \text{ with }\underline 0:=\mathsf{tree}\,\mathbf 0\,(),\ \underline{n+1}:=\underline n\cup\{\underline n\}\\
\{a,b\}^* &:= \mathsf{tree}\,\mathbf 2\,(\mathsf{ind}^{\mathbf 2}\,a\,b)\qquad
(\textstyle\bigcup a)^* := \mathsf{tree}\,(\Sigma i{:}\bar a.\,\overline{f_a i})\,(\langle i,j\rangle\mapsto f_{f_a i}\,j)\\
\PP(a)^* &:= \mathsf{tree}\,(\bar a\to\Prop)\,\bigl(P\mapsto\mathsf{tree}\,(\Sigma i{:}\bar a.\,P\,i)\,(f_a\circ\pi_1)\bigr)\\
\{x\in a\mid\varphi\}^* &:= \mathsf{tree}\,(\Sigma i{:}\bar a.\,\varphi^*(f_a\,i))\,(f_a\circ\pi_1)\qquad
\{t\mid x\in u\}^* := \mathsf{tree}\,\bar u\,(i\mapsto t^*(f_u\,i))\\
\rec_{x,g.G}(t)^* &:= r(t^*)\ \text{ where } r:=\mathsf{ind}^W_{\mathsf{set}}\bigl(\lambda X\,f\,r'.\ G^*(\mathsf{tree}\,X\,f,\ \mathsf{tree}\,X\,(i\mapsto\langle f\,i,\,r'\,i\rangle))\bigr)
\end{align*}
and $(t=u)^*:=t^*\approx u^*$, $(t\in u)^*:=t^*\in u^*$, connectives
pointwise, $(\forall x\,\varphi)^*:=\Pi x{:}\mathsf{set}.\,\varphi^*$,
$(\exists x\,\varphi)^*:=\exists x{:}\mathsf{set}.\,\varphi^*$. The
$\wfrec$ clause is given in Lemma~\ref{lem:wfrec}.
\end{definition}

The first six clauses are Werner's and Barras's; the image clause is
Barras's axiom-free $\mathsf{replf}$ \cite{barras10}. Note where the
universes are used: $\bar a\to\Prop$ and $\Sigma i{:}\bar a.\,P\,i$ are in
$\Type_0$ because $\Prop:\Type_0$ and $\bar a:\Type_0$, so $\PP(a)^*$ and
comprehension are trees, and the quantifiers in $\varphi^*$ range over
$\mathsf{set}:\Type_1$ inside $\Prop$ by impredicativity. The recursion
clause is $W$-elimination with motive $\mathsf{set}$, a type in $\Type_1$,
which is allowed since motives may be in any sort.

\begin{lemma}[compatibility]\label{lem:leibniz}
Every $t^*$ respects $\approx$ in each argument, and every $\varphi^*$
is invariant under $\approx$.
\end{lemma}
\begin{proof}
Induction on syntax. The base formers and the propositional clauses are
\cite[Prop.~2]{miquel06} transported to trees, using that $\in$ respects
$\approx$. For $\rec$: show $a\approx a'\to r(a)\approx r(a')$ by
$W$-induction on $a$. Writing $a=\mathsf{tree}\,X\,f$ and
$a'=\mathsf{tree}\,X'\,f'$, from $a\approx a'$ every $f\,i$ is $\approx$ to
some $f'\,j$ and conversely, so by the induction hypothesis the
restriction trees $\mathsf{tree}\,X\,(i\mapsto\langle f\,i,r\,(f\,i)\rangle)$
and $\mathsf{tree}\,X'\,(j\mapsto\langle f'\,j,r\,(f'\,j)\rangle)$ are
bisimilar, and $G^*$ respects $\approx$ by the outer induction.
\end{proof}

\begin{lemma}[well-foundedness gives accessibility]\label{lem:wf}
Let $X,R:\mathsf{set}$ and define $R'\,i\,j:=\langle f_X\,i,f_X\,j\rangle\in R$
on $\bar X$. Then $\CIC_2+\EM\vdash\WF(X,R)^*\to\Pi i{:}\bar X.\,\mathsf{Acc}_{R'}\,i$.
\end{lemma}
\begin{proof}
Assume $\WF(X,R)^*$ and, towards a contradiction, $\neg\mathsf{Acc}_{R'}\,i_0$.
Let $S:=\mathsf{tree}\,(\Sigma j{:}\bar X.\,\neg\mathsf{Acc}_{R'}\,j)\,(f_X\circ\pi_1)$,
a set with $S\subseteq X$ and $f_X\,i_0\in S$. By $\WF$ there is an
$R$-minimal $m\in S$, so $m\approx f_X\,j$ with $\neg\mathsf{Acc}_{R'}\,j$
and no element of $S$ is $R$-below $m$. For any $k$ with $R'\,k\,j$ we get
$f_X\,k\notin S$, hence $\neg\neg\mathsf{Acc}_{R'}\,k$, hence
$\mathsf{Acc}_{R'}\,k$ by $\mathsf{em}$. So $\mathsf{acc}\,j\,(\ldots)$
proves $\mathsf{Acc}_{R'}\,j$, a contradiction; by $\mathsf{em}$ again,
$\mathsf{Acc}_{R'}\,i_0$.
\end{proof}

\begin{lemma}[$\wfrec$]\label{lem:wfrec}
There is a term $\wfrec_{x,g.G}(X,R,t)^*$ such that $\CIC_2+\EM$ proves
(WfRec)$^*$ and Lemma~\ref{lem:leibniz} for it.
\end{lemma}
\begin{proof}
Given $X,R$, first define
$r:\Pi i{:}\bar X.\,\mathsf{Acc}_{R'}\,i\to\mathsf{set}$ by
$\mathsf{ind}^{\mathsf{Acc}}$ with motive $\mathsf{set}$:
\[
r\,i\,(\mathsf{acc}\,i\,h):=G^*\bigl(f_X\,i,\ \mathsf{tree}\,(\Sigma j{:}\bar X.\,R'\,j\,i)\,(\langle j,\rho\rangle\mapsto\langle f_X\,j,\ r\,j\,(h\,j\,\rho)\rangle)\bigr).
\]
By $\mathsf{Acc}$-induction and proof irrelevance the value does not
depend on the accessibility proof, and, using Lemma~\ref{lem:leibniz}
for $G$, $f_X\,i\approx f_X\,i'$ implies $r\,i\approx r\,i'$ (the sets of
$R'$-predecessors correspond under $\approx$). Now put
\[
\wfrec_{x,g.G}(X,R,a)^*:=\textstyle\bigcup\mathrm{Vals}_a,\qquad
\mathrm{Vals}_a:=\mathsf{tree}\,\bigl(\Sigma i{:}\bar X.\,(f_X\,i\approx a)\wedge\mathsf{Acc}_{R'}\,i\bigr)\,\bigl(\langle i,p\rangle\mapsto r\,i\,(\pi_2 p)\bigr),
\]
the union of the set of values at all indices representing $a$; the index
type is in $\Type_0$ since its components are propositions. If
$\WF(X,R)^*$ and $a\in X$, then by
Lemma~\ref{lem:wf} every $i$ is accessible, some $i$ has $f_X\,i\approx a$,
all such $i$ give $\approx$-equal $r\,i$, and the union of a nonempty
family of pairwise $\approx$-equal sets is $\approx$ to any member. So
the definition avoids choosing an index: it is the standard set-theoretic
device that the unique member of a nonempty subsingleton family is its
union, applied to a family indexed by the \emph{type} $\bar X$, which is
what makes the family a set. (WfRec)$^*$ then follows from the defining
equation of $r$, and compatibility in $X$, $R$, $a$ from
Lemma~\ref{lem:leibniz} and $\mathsf{Acc}$-induction. If the hypothesis
fails the term still denotes some set, as (WfRec) requires nothing.
\end{proof}

\begin{theorem}\label{thm:interp}
For every sentence $\varphi$ of $\Lang$, if $\ZFh\vdash\varphi$ then
$\CIC_2+\EM\vdash\varphi^*$. In particular $\CIC_2+\EM$ interprets
$\ZFh$.
\end{theorem}
\begin{proof}
Classical first-order logic with equality is sound for the translation:
$\mathsf{em}$ gives classical logic in $\Prop$, $\approx$ is an
equivalence relation, and Lemma~\ref{lem:leibniz} gives the congruence
axioms for every function and relation symbol. The axioms: (Ext) is
the definition of $\approx$ and $\in$ (the tree-to-tree direction uses
the two halves of $\approx$). (Pair), (Union), (Pow), (Inf), (Comp) are
\cite{werner97,barras10}; for (Pow) one needs, given $b\subseteq a$, a
$P:\bar a\to\Prop$ with $\mathsf{tree}(\Sigma i.\,P\,i)(f_a\circ\pi_1)\approx b$,
namely $P\,i:=f_a\,i\in b$. (Img) is $\mathsf{replf}$: $y\in\{t\mid x\in u\}^*$
iff $\exists i{:}\bar u.\,y\approx t^*(f_u\,i)$ iff $\exists x\,(x\in u\wedge y=t)^*$,
using Lemma~\ref{lem:leibniz} for the right-to-left direction. (Rec) is
the computation rule of $\mathsf{ind}^W$ together with the observation
that $\{\langle y,\rec(y)\rangle\mid y\in a\}^*$ is exactly the restriction
tree in Definition~\ref{def:star}. (WfRec) is Lemma~\ref{lem:wfrec}.
($\in$-Ind) is $W$-induction: from $\forall x\,(\forall y{\in}x\,\varphi\to\varphi(x))^*$
prove $\Pi a{:}\mathsf{set}.\,\varphi^*(a)$ by $\mathsf{ind}^W$ with
motive $\varphi^*(a):\Prop$, since $y\in\mathsf{tree}\,X\,f$ means
$y\approx f\,i$ for some $i$ and $\varphi^*$ is $\approx$-invariant.
\end{proof}

\begin{proposition}\label{prop:con}
$\CIC_2+\EM\vdash\Con(\ZFh)$. Hence, if $\CIC_2+\EM$ is consistent, no
interpretation of $\CIC_2+\EM$ in $\ZFh$ preserves arithmetic.
\end{proposition}
\begin{proof}[Proof sketch]
The syntax of $\Lang$ is a $W$-type over $\N$ in $\Type_0$. Define by
$W$-recursion an evaluation
$\mathsf{ev}:\Terms\to(\mathrm{Var}\to\mathsf{set})\to\mathsf{set}$, with
motive in $\Type_1$, following Definition~\ref{def:star}, and a
satisfaction predicate $\mathsf{Sat}:\mathrm{Form}\to(\mathrm{Var}\to\mathsf{set})\to\Prop$
with $\mathsf{Sat}(\forall x\,\varphi)\rho:=\Pi a{:}\mathsf{set}.\,\mathsf{Sat}(\varphi)(\rho[x{:=}a])$,
which is a proposition by impredicativity. The proof of
Theorem~\ref{thm:interp} then internalises: by induction on codes of
derivations, every theorem of $\ZFh$ is $\mathsf{Sat}$-true, and $\bot$ is
not. This is the usual argument that a theory with a truth predicate for
$T$ over a class-sized domain proves $\Con(T)$; the ingredient specific
to type theory is that $\mathsf{set}$ is a type and $\Prop$ is
impredicative, so quantification over all sets inside a proposition is
available. The second claim follows since an arithmetic-preserving
interpretation would let $\ZFh$ prove $\Con(\ZFh)$.
\end{proof}

So the internal set theory of $\CIC_2+\EM$ is $\ZFh$, but the type theory
is strictly stronger than $\ZFh$: it has, in addition, an impredicative
theory of classes over $\mathsf{set}$, in the way Kelley--Morse is
stronger than $\ZF$. It does not, however, reach $\ZF$: Replacement
fails for functional relations because $\exists$ is truncated, as
recorded in \cite[Remark~2.1, Prop.~5.1]{paper1}. Adding $\uchoice$ gives
$\Con(\ZF)$ \cite{barras10,friedman73}.

\section{Universes}\label{sec:univ}

\subsection{$\ZFh_n$ and $\ZFh_\omega$}

Let $\Lang_n$ be $\Lang$ with new constants $U_0,\dots,U_{n-1}$, and
$\Lang_\omega$ with constants $U_i$ for all $i<\omega$. Terms and
formulas of $\Lang_n$ may mention the constants anywhere.

\begin{definition}
$\ZFh_n$ is $\ZFh$ formulated in $\Lang_n$ (all schemata extended to
$\Lang_n$) plus, for each $i<n$:
\begin{itemize}
\item $U_i$ is transitive, and $U_i\in U_{i+1}$ when $i+1<n$;
\item (Closure) for every term $t(\vec p)$ of $\Lang_n$ whose constants
  are among $U_0,\dots,U_{i-1}$:\ $\forall\vec p\in U_i.\ t(\vec p)\in U_i$.
\end{itemize}
$\ZFh_\omega$ is the union of the $\ZFh_n$ over the common language $\Lang_\omega$.
\end{definition}

Closure includes powerset, comprehension for arbitrary formulas (whose
quantifiers range over the whole universe), images, and both recursions.
So each $U_i$ is a transitive model of $\ZFh$ with Separation for
formulas with unbounded quantifiers, and in particular
$\ZFh_{n+1}\vdash\Con(\ZFh_n)$ by the satisfaction relation for the set
$U_n$.

\begin{theorem}\label{thm:univ}
$\ZF\vdash\Con(\ZFh_\omega)$. More precisely, if $\lambda_0<\lambda_1<\cdots$
are $\Lang$-closed ordinals with each $\lambda_{i+1}$ closed under the
majorants of $\Lang_{i+1}$-terms in the sense below, and
$\lambda_\omega:=\sup_i\lambda_i$, then $V_{\lambda_\omega}$ with
$U_i:=V_{\lambda_i}$ satisfies $\ZFh_\omega$, and $\ZF$ proves such
ordinals exist.
\end{theorem}
\begin{proof}
Extend Definition~\ref{def:maj} by $f_{U_i}(\alpha)=\lambda_i$; with
this, Lemma~\ref{lem:rank} holds verbatim for $\Lang_\omega$-terms
evaluated with $\sem{U_i}=V_{\lambda_i}$. Choose $\lambda_i$ by
Lemma~\ref{lem:closed} successively, taking $\alpha_0>\lambda_{i-1}$
and closing under the majorants of all $\Lang_i$-terms, which is
legitimate because $f_{U_j}$ for $j<i$ is the constant $\lambda_j$,
already fixed. Then $V_{\lambda_\omega}$ is a limit stage,
$\lambda_\omega$ is $\Lang_\omega$-closed since every $\alpha<\lambda_\omega$
is below some $\lambda_i$, and Theorem~\ref{thm:main} applies to give
the $\ZFh$ axioms. $V_{\lambda_i}$ is transitive and
$V_{\lambda_i}\in V_{\lambda_{i+1}}$. Closure: for a term $t$ with
constants among $U_j$, $j<i$, and $\vec p\in V_{\lambda_i}$, the rank
lemma gives $\rk\sem t\le f_t(\alpha)$ for $\alpha<\lambda_i$, and
$\lambda_i$ is closed under $f_t$, so $\sem t\in V_{\lambda_i}$.
\end{proof}

\subsection{Relation to $\CIC_\omega+\EM$}

Let $\CIC_{n+2}+\EM$ have sorts $\Prop:\Type_0:\cdots:\Type_{n+1}$ with
$\Type_{n+1}$ the top sort, and $\CIC_\omega+\EM$ the union. Put
$\mathsf{set}_i:=Wx{:}\Type_i.\,x$, a type in $\Type_{i+1}$; by
cumulativity $\mathsf{set}_i\subseteq\mathsf{set}_{i+1}$ literally, and
$\mathsf{set}:=\mathsf{set}_n$ is the ambient set type of
$\CIC_{n+2}+\EM$. The universe constants are interpreted by
\[
U_i^*:=\mathsf{tree}\ \mathsf{set}_i\ (a\mapsto a)\qquad(i<n),
\]
which is a tree over the type $\mathsf{set}_i:\Type_{i+1}\subseteq\Type_n$,
hence an element of $\mathsf{set}$, and an element of $U_{i+1}^*$ as
$\mathsf{set}_i:\Type_{i+1}$.

\begin{theorem}\label{thm:interpn}
$\CIC_{n+2}+\EM$ interprets $\ZFh_n$, and $\CIC_\omega+\EM$ interprets
$\ZFh_\omega$. Moreover $\CIC_{n+2}+\EM\vdash\Con(\ZFh_n)$.
\end{theorem}
\begin{proof}
Theorem~\ref{thm:interp} with $\Type_0,\Type_1$ replaced by
$\Type_n,\Type_{n+1}$ gives the $\ZFh$ axioms over $\mathsf{set}$. $U_i^*$
is transitive because an element of $U_i^*$ is $\approx$ to a tree in
$\mathsf{set}_i$, whose subtrees are in $\mathsf{set}_i$. Closure: by
induction on the term $t$, if all parameters are in $\mathsf{set}_i$
then $t^*(\vec p)$ is a tree all of whose index types are in $\Type_i$,
since each clause of Definition~\ref{def:star} forms index types by
$\Sigma$, $\to\Prop$, and finite types from index types of its
arguments, and $\rec$, $\wfrec$ produce values by recursion whose every
step does the same; so $t^*(\vec p)\in\mathsf{set}_i$ up to $\approx$,
and $U_i^*$ contains it. (For terms mentioning $U_j$, $j<i$: $U_j^*$
itself has index type $\mathsf{set}_j:\Type_{j+1}\subseteq\Type_i$.)
$\Con(\ZFh_n)$ is Proposition~\ref{prop:con} over $\mathsf{set}_n$.
\end{proof}

\paragraph{The picture.}
For each $n$ we have, in consistency strength and provably in $\ZF$,
\[
\ZFh_n\ <\ \CIC_{n+2}+\EM\ \le^{?}\ \ZFh_{n+1}\ <\ \CIC_{n+3}+\EM\ \le^{?}\ \cdots\ <\ \ZF,
\]
where $<$ is ``proves the consistency of'' (Theorem~\ref{thm:interpn}
and $\ZFh_{n+1}\vdash\Con(\ZFh_n)$), the final $<$ is Theorem~\ref{thm:univ},
and $\le^{?}$ is the conjectured upper bound: a model of $\CIC_{n+2}+\EM$
in $\ZFh_{n+1}$, interpreting $\Type_i$ as $U_i$ for $i\le n$ and
$\Type_{n+1}$ as the universe, with function types interpreted by
term-definable functions so that the majorant argument of
Section~\ref{sec:rank}, applied to the type theory's own terms, bounds
every $\Sigma$ over a family. The half-step from $\ZFh_n$ to
$\CIC_{n+2}+\EM$ is an impredicative class theory over $\mathsf{set}_n$;
the half-step from $\CIC_{n+2}+\EM$ to $\ZFh_{n+1}$ is that a class
theory over $U_n$ becomes a set theory once $U_n$ is a set. Both
half-steps are far below an inaccessible: everything in the tower is
modelled in $V_{\lambda_\omega}$ for the countable-cofinality
$\lambda_\omega$ of Theorem~\ref{thm:univ}.

\section{The upper bound: hereditarily bounded set models}\label{sec:upper}

This section sets out the model that would give the $\le^{?}$ steps, and
proves the two lemmas that address the obstruction met in
\cite[Prop.~4.4]{paper1}. What it does not do is complete the soundness
proof; the status is recorded at the end. The reader should know in
advance that the closure ordinals as defined in \S\ref{sec:closure} are
too small (Example~\ref{ex:sat}), so that the model of \S\ref{sec:adm} is
not a model; Remark~\ref{rem:nosat} gives the repair we expect to work,
under which most of \S\ref{sec:adm}--\ref{sec:transfer} becomes
unnecessary.

\paragraph{Why a set model, and which one.}
The natural model of $\CIC_\omega+\EM$ is Werner's: $\Prop$ is
$\{\bot,\top\}$, types are sets, $\Pi$ is the set-theoretic dependent
product, $\Type_i$ is $V_{\kappa_i}$. It needs $\kappa_i$ inaccessible
for exactly one reason: in the context $X{:}\Type_0,\ f{:}X\to\Type_0$,
the semantic value of $f$ ranges over \emph{all} functions
$|X|\to V_{\kappa_0}$, and $\Sigma x{:}X.\,f\,x$ must land in
$V_{\kappa_0}$. Replacing set-theoretic function spaces by Kleene-tracked
ones does not help \cite[Prop.~4.4]{paper1}, and term models do not
survive an impredicative $\Prop$ with excluded middle. So the function
spaces must be cut down by a semantic condition that (a) every
term-definable function satisfies, (b) is preserved by application,
recursion, and abstraction, and (c) bounds $\sup_x\rk f(x)$ below
$\lambda_0$ whenever the domain is small. Condition (c) is the content of
Section~\ref{sec:rank} for $\Lang$-terms: a definable family has a rank
majorant computed from its shape. Conditions (a)--(b) say the majorant
must be attached to \emph{elements}, not just to terms, so that a
variable $f$ in a context carries one. This is Howard's majorizability,
with ordinals in place of numbers and with functionals over ordinals as
majorants at higher types.

\paragraph{Plan of the construction.}
The definitions below are given in the order in which they are used,
and the construction has two phases. The first phase determines the
universe extents $\lambda_0<\lambda_1<\cdots$ by a recursion on the
level $N$: codes and their full decoding (\S\ref{sec:codes}) and the
full evaluation of terms with its hybrid majorants (\S\ref{sec:hybrid})
are defined once and for all, independently of the extents except for
the universe constants; then
(\S\ref{sec:closure}) stage $N$ of the recursion, given
$\lambda_0,\dots,\lambda_{N-1}$, defines closure functions
$\nu^N_{t,\theta}$ by taking suprema of majorants over auxiliary
structures whose universe $U_N$ is cut at a variable ordinal $\alpha$
and whose higher universes are arbitrary, and sets $\lambda_N$ to be the
least ordinal closed under them. Nothing in this phase refers to
$\lambda_j$ for $j\ge N$, and nothing in it refers to the model. The
second phase (\S\ref{sec:adm}--\ref{sec:transfer}) takes all the
$\lambda_j$ as fixed and defines the model: the retraction
$r^N_\alpha$ from the model onto the auxiliary structure with $U_N$ cut
at $\alpha$ and the true higher universes, the admissible elements, and
the decoding $\El$. Its main lemma, the transfer, says that the
retraction of an admissible environment is one of the environments over
which $\nu^N$ took its supremum; this is the sense in which the first
phase chose the $\lambda_N$ correctly, and it is why the suprema of the
first phase range over arbitrary higher extents. Forward references in
the first phase to the model are motivation only.

\subsection{Codes and full decoding}\label{sec:codes}

Throughout, $\lambda_0<\lambda_1<\cdots$ are the extents produced by
\S\ref{sec:closure}, with supremum $\lambda_\omega$; until that section
they are arbitrary ordinals, and the definitions of this section and
the next do not depend on them except through the universes
$U_i:=\{c:\rk(c)<\lambda_i\}$.

\begin{definition}[codes]\label{def:codes2}
The class of \emph{codes} is defined by recursion on rank: the constants
$\bot,\top,\code\Prop,\code{\mathbf n},\code\N$ are codes, of finite
rank, where $\bot$ is the empty set; the \emph{universe constant} $\code{\U_i}:=\langle\lambda_i,\mathsf U\rangle$
is a code, of rank $\lambda_i+1$; if $c$ is a code, $G$ a function from
$V_{\rk(c)\cdot\omega}$ to codes, and $R\subseteq V_{\rk(c)\cdot\omega}^2$,
then $\code\Pi(c,G)$, $\code\Sigma(c,G)$, $\code W(c,G)$, $\code/(c,R)$
are codes. Codes of rank below any $\alpha$ form a set. Put
$U_i:=\{c:\rk(c)<\lambda_i\}$ and $\ell(c):=$ the least $i$ with $c\in U_i$,
the \emph{sort level}; so $\code{\U_i}\in U_{i+1}\setminus U_i$. The rank
of the universe constant is chosen so that $U_i\subseteq V_{\rk(\code{\U_i})\cdot\omega}$,
which the family domains below require. Consequently the codes that
mention no universe constant are independent of the $\lambda_j$, and a
code mentioning $\code{\U_j}$ depends on $\lambda_j$ only through this
rank; in the auxiliary structures of \S\ref{sec:closure}, where $U_j$
is cut at some ordinal $\gamma$, the universe constant is
$\langle\gamma,\mathsf U\rangle$. The \emph{value level} is
$\kappa(\code{\U_i}):=i$, $\kappa:=0$ for the other constants,
$\kappa(\code\Pi(c,G)):=\sup_a\kappa(G(a))$,
$\kappa(\code\Sigma(c,G))=\kappa(\code W(c,G)):=\max(\kappa(c),\sup_a\kappa(G(a)))$,
$\kappa(\code/(c,R)):=\kappa(c)$.
\end{definition}

\begin{definition}[majorants]\label{def:M}
For an ordinal $\theta$ define sets $\mathbb M_\theta^{(\delta)}$ by
recursion on $\delta$. A \emph{majorant of degree $<\delta$ over $\theta$} is a
triple $m=(m_{\mathrm o},m_{\mathrm p},m_{\mathrm f})$ of an \emph{ordinal
part} $m_{\mathrm o}\in\theta$, a \emph{pair part} $m_{\mathrm p}=\langle m_1,m_2\rangle$
with $m_1,m_2\in\mathbb M_\theta^{(\beta)}$, and a \emph{function part}
$m_{\mathrm f}:\mathbb M_{\theta'}^{(\beta)}\to\mathbb M_\theta^{(\beta)}$ monotone,
for some $\beta<\delta$ and $\theta'\le\theta$ (the domain may be
over a smaller ordinal than the values, as for $\lambda x{:}\N.\,\Type_0$),
each of the three parts being optional
(written $-$ when absent), with at least one present.
$\mathbb M_\theta:=\bigcup_\delta\mathbb M_\theta^{(\delta)}$ is a set, and
$\mathbb M:=\bigcup_\theta\mathbb M_\theta$ is the class of majorants; each
majorant is a set. The order is $m\le m'$ iff every part present in $m$
is present in $m'$ and is $\le$ there: ordinals as ordinals, pairs
componentwise, functions pointwise on the common domain. Joins
$m\sqcup m'$ and suprema of sets of majorants are taken componentwise,
an absent part being neutral. A function part is applied to an
arbitrary majorant $\mu$, inside or outside its domain, by monotone
extension, $m_{\mathrm f}(\mu):=\sup\{m_{\mathrm f}(\chi):\chi\in\mathrm{dom}\,m_{\mathrm f},\ \chi\le\mu\}$,
a set supremum; on the domain this is $m_{\mathrm f}$ itself, and every
value of the extension is a supremum of values on the domain. The \emph{leaves} of $m$ are the ordinals
occurring in it, recursively through pairs and through all values of
function parts; $m$ is \emph{hereditarily below $\lambda$} if all its
leaves are, and the \emph{leaf supremum} of $m$ is their supremum. We
write $\beta$ for the majorant $(\beta,-,-)$, and $\hat\beta$ for the
\emph{hereditary constant}, whose ordinal part is $\beta$, whose pair
part is $\langle\hat\beta,\hat\beta\rangle$, and whose function part is
constantly $\hat\beta$, at each degree.
\end{definition}

The three parts exist because the shape of a fibre varies with the
argument of a dependent function: for $f:\Pi b{:}\mathbf 2.\,B(b)$ with
$B(0)$ a base type and $B(1)$ a function type, the canonical majorant of
$f$ at an input covering both $0$ and $1$ must bound an ordinal-shaped
element and a function at once (\S\ref{sec:examples}). Keeping the parts
separate and joins componentwise is essential: reading the ordinal part
of a big code into the function part would let a $\Type_2$ branch bound
the values of a $\Type_0\to\Type_0$ branch, and the pollution of
Example~\ref{ex:hetero} would return through the join.

\begin{definition}[full decoding]\label{def:Elf}
By recursion on codes define the set $\El^{\mathrm f}(c)$, Werner's decoding:
\begin{align*}
\El^{\mathrm f}(\bot)&:=\emptyset,\quad \El^{\mathrm f}(\top):=\{\star\},\quad
\El^{\mathrm f}(\code\Prop):=\{\bot,\top\},\\
\El^{\mathrm f}(\code{\mathbf n})&:=n,\quad
\El^{\mathrm f}(\code\N):=\omega,\quad \El^{\mathrm f}(\code{\U_i}):=U_i;\\
\El^{\mathrm f}(\code\Pi(c,G))&:=\{f:\ f\text{ a function on }\El^{\mathrm f}(c)\text{ with }f(a)\in\El^{\mathrm f}(G(a))\};\\
\El^{\mathrm f}(\code\Sigma(c,G))&:=\{\langle a,b\rangle:\ a\in\El^{\mathrm f}(c),\ b\in\El^{\mathrm f}(G(a))\};\\
\El^{\mathrm f}(\code W(c,G))&:=\text{well-founded trees, as sets of finite paths, with labels }a\in\El^{\mathrm f}(c)\\
&\phantom{:=}\text{ and children indexed by }\El^{\mathrm f}(G(a));\\
\El^{\mathrm f}(\code/(c,R))&:=\El^{\mathrm f}(c)/{\sim_R},\ \text{the classes of the equivalence relation generated by }R;
\end{align*}
with propositional codes identified with truth values as in
\cite[Def.~4.1]{paper1}. By induction on codes,
$\El^{\mathrm f}(c)\subseteq V_{\rk(c)\cdot\omega}$, so every family $G$ is
defined on all of $\El^{\mathrm f}(c)$.
\end{definition}

\begin{definition}[majorization and canonical majorants]\label{def:codesbd}
By recursion on codes define, for $e\in\El^{\mathrm f}(c)$ and $m\in\mathbb M$,
the relation $e\lhd_c m$ (``$m$ majorizes $e$'') and the canonical
majorant $\hat e\in\mathbb M$:
\begin{itemize}
\item Constants and universes: $e\lhd m$ iff $m_{\mathrm o}$ is present and
  $\rk e<m_{\mathrm o}$; $\hat e:=\rk e+1$.
\item Pairs, $c=\code\Sigma(c',G)$: $\langle a,b\rangle\lhd m$ iff
  $m_{\mathrm p}=\langle m_1,m_2\rangle$ is present, $a\lhd_{c'}m_1$, and
  $b\lhd_{G(a)}m_2$; $\widehat{\langle a,b\rangle}:=(-,\langle\hat a,\hat b\rangle,-)$.
\item Trees, $c=\code W(c',G)$: $t\lhd m$ iff $m_{\mathrm p}=\langle h,m'\rangle$
  is present, $h$ is an ordinal exceeding the height of $t$, and every
  label of $t$ is majorized by $m'$; $\hat t:=(-,\langle\mathrm{height}(t)+1,\ \sup_a\hat a\rangle,-)$
  over the labels $a$ of $t$.
\item Classes, $c=\code/(c',R)$: $[x]\lhd m$ iff $x'\lhd_{c'}m$ for some
  $x'\in[x]$; $\widehat{[x]}:=$ the pointwise infimum of the $\hat{x}'$
  over the members.
\item Functions, $c=\code\Pi(c',G)$: $f\lhd m$ iff $m_{\mathrm f}$ is present
  with domain containing $\mathbb M_{\theta}^{(\rk c'+1)}$ for
  $\theta:=\rk(c')\cdot\omega$, and $f(a)\lhd_{G(a)}m_{\mathrm f}(\hat a)$ for
  every $a\in\El^{\mathrm f}(c')$; and
  $\hat f:=(-,-,\varphi)$ with $\varphi(\chi):=\sup\{\widehat{f(a)}:a\in\El^{\mathrm f}(c'),\ \hat a\le\chi\}$
  on that domain.
\item Truth values: $\star\lhd m$ iff $m_{\mathrm o}$ is present and
  positive; $\hat\star:=1$.
\end{itemize}
\end{definition}

\begin{lemma}\label{lem:Mbasic}
For $e\in\El^{\mathrm f}(c)$: (1) $\hat e\in\mathbb M_{\rk(c)\cdot\omega}^{(\rk c+1)}$
and $e\lhd\hat e$; (2) $\hat e$ is the least majorant of $e$, and
$e\lhd m\le m'$ implies $e\lhd m'$; (3) $\hat f$ is monotone, and
$\hat e$ is hereditarily below $\rk(c)\cdot\omega$; (4) the suprema and
joins in Definition~\ref{def:M} preserve monotonicity of function parts.
\end{lemma}
\begin{proof}
All by induction on codes. For (1) at $\Pi$, the domain requirement is
met because $\hat a\in\mathbb M^{(\rk c'+1)}_{\rk(c')\cdot\omega}$ by induction,
and the values $\widehat{f(a)}$ have degree $\le\rk G(a)+1<\rk c+1$. For
(2) at $\Pi$, $f\lhd m$ gives $\widehat{f(a)}\le m_{\mathrm f}(\hat a)\le m_{\mathrm f}(\chi)$
for $\hat a\le\chi$ by monotonicity, so $\varphi\le m_{\mathrm f}$; upward
closure uses monotonicity of $m'_{\mathrm f}$ in the same way. (3) is
$\El^{\mathrm f}(c)\subseteq V_{\rk(c)\cdot\omega}$.
\end{proof}

\begin{example}[why the $\lambda_i$ must be closed under Hartogs numbers]\label{ex:set}
Let $\mathsf{set}:=Wx{:}\Type_0.\,x$ and define $T:\mathsf{set}\to\Type_0$ by
$W$-recursion, $T(\mathsf{tree}\,X\,d):=(\Sigma x{:}X.\,T(d\,x))\to\Prop$.
Then $\rk T(t)\ge\omega+\mathrm{height}(t)$, and $T(t)$ must lie in $U_0$,
so every tree over a small type must have height below $\lambda_0$. A
tree over $X\in U_0$ has height below $\hartogs(\PP(X))$, and trees over
$\N$ realise every countable height, so $\lambda_0$ must be closed under
$\alpha\mapsto\hartogs(V_{\alpha+1})$; in $\ZFC$ this makes it a strong
limit cardinal, of cofinality $\omega$ if obtained by an $\omega$-iteration.
This is the ordinal Miquel's footnote and Section~\ref{sec:rank} both
point at.
\end{example}

The model will consist of the \emph{admissible} elements of the full
decoding, with function spaces restricted to admissible arguments. Which
elements are admissible is decided by term majorants, defined next
without reference to admissibility.

\subsection{Evaluation and hybrid term majorants}\label{sec:hybrid}

\begin{definition}[$N$-small and $N$-big variables]\label{def:smallbig}
Let $\Gamma$ be a context, $x{:}A\in\Gamma$, and $N\ge1$. The variable
$x$ is \emph{$N$-small} if $\Gamma'\vdash A:\Prop$ or $\Gamma'\vdash A:\Type_j$
with $j<N$ is derivable, where $\Gamma'$ is the prefix of $\Gamma$ before
$x$, and \emph{$N$-big} otherwise, i.e.\ if the least sort of $A$ is
$\Type_j$ with $j\ge N$. Since sorts are cumulative and every type has
a least sort, this is a syntactic property of the context. ``Small'' and
``big'' without qualification mean $1$-small and $1$-big. Semantically,
by Lemma~\ref{lem:tyform}, an $N$-small variable's type decodes to a code
of level below $N$, whose full decoding is fixed by every retraction
$r^N_\alpha$; an $N$-big variable ranges over elements that the level-$N$
construction must control.

The \emph{syntactic value level} $\vl(T)$ of a type $T$ is defined by
$\vl(\Prop):=0$, $\vl(\Type_i):=i$, $\vl(\Pi x{:}A.B):=\vl(B)$,
$\vl(\Sigma x{:}A.B)=\vl(Wx{:}A.B):=\max(\vl A,\vl B)$, $\vl(A/R):=\vl(A)$,
and $\vl(T):=j$ for any other type $T$ whose least sort is $\Type_j$
(and $0$ if it is $\Prop$). The value level of a term $t$ with
$\Gamma\vdash t:T$ is $\vl(T)$; it bounds the value level $\kappa$ of
the code of $T$ in every environment, and it is the level at which the
closure functions of \S\ref{sec:closure} will bound $t$.
\end{definition}

\begin{definition}[full evaluation]\label{def:fulleval}
For a term $t$ and an environment $\rho$ assigning sets to variables,
define $\sem t\rho$ by recursion on $t$; the clauses are Werner's, made
total by returning $\emptyset$ where a clause does not apply. Since
$\emptyset$ is the code $\bot$, a family $a\mapsto\sem B\rho[x{:=}a]$ is
code-valued on its whole domain $V_{\rk\sem A\rho\cdot\omega}$ whenever
$B$ is a type, and so every type former evaluates to a code in every
environment, junk or not.
\begin{itemize}
\item $\sem x\rho:=\rho(x)$; $\sem{\Prop}\rho:=\code\Prop$, $\sem{\Type_i}\rho:=\code{\U_i}$,
  $\sem{\mathbf n}\rho:=\code{\mathbf n}$, $\sem\N\rho:=\code\N$.
\item $\sem{\Pi x{:}A.B}\rho:=\code\Pi(\sem A\rho,\ a\mapsto\sem B\rho[x{:=}a])$
  with the family taken on $V_{\rk\sem A\rho\cdot\omega}$, and likewise
  $\Sigma$ and $W$; $\sem{A/R}\rho:=\code/(\sem A\rho,\{(a,a'):\sem R\rho[x{:=}a,x'{:=}a']=\top\})$;
  propositional formers are identified with truth values as in
  \cite[Def.~4.1]{paper1}, with $\sem{a=_A a'}\rho:=[\sem a\rho=\sem{a'}\rho]$,
  $\sem{\|A\|}\rho:=[\El^{\mathrm f}(\sem A\rho)\ne\emptyset]$, and
  $\sem{\Pi x{:}A.P}\rho:=[\forall a\in\El^{\mathrm f}(\sem A\rho).\ \sem P\rho[x{:=}a]=\top]$
  when $P$ is a proposition.
\item $\sem{\lambda x.b}\rho:=$ the function $a\mapsto\sem b\rho[x{:=}a]$ on
  $\El^{\mathrm f}(\sem A\rho)$; $\sem{f\,a}\rho:=\sem f\rho(\sem a\rho)$;
  $\sem{\langle a,b\rangle}\rho:=\langle\sem a\rho,\sem b\rho\rangle$;
  $\sem{\mathsf{tree}\,a\,d}\rho:=$ the tree with root $\sem a\rho$ and
  subtrees $\sem d\rho(y)$; $\sem{[a]_R}\rho:=$ the class of $\sem a\rho$;
  $\sem{\mathbf k_{\mathbf n}}\rho:=k$, $\sem 0\rho:=0$, $\sem{\mathsf S\,m}\rho:=\sem m\rho+1$;
  all proofs and propositional constants evaluate to $\star$.
\item The eliminators are the corresponding set-theoretic recursions on
  the actual value of their argument: $\mathsf{ind}^{\mathbf n}$ and
  $\mathsf{ind}^\Sigma$ by case, $\mathsf{ind}^\N$ by primitive recursion,
  $\mathsf{ind}^W$ by recursion on the well-founded tree,
  $\mathsf{ind}^{\mathsf{Acc}}$ by recursion along the accessible part of
  the set relation $\sem R\rho$, $\mathsf{ind}^=$ by $f$ at the argument,
  $\mathsf{ind}^{\|\|}$ to $\star$, and $\mathsf{ind}^/f\,h\,q$ to the common
  value of $f$ on the members of $q$ if they agree and $\emptyset$
  otherwise.
\end{itemize}
The \emph{model's} interpretation, defined in \S\ref{sec:adm}, is this
evaluation with abstractions restricted to admissible arguments; on
admissible environments the two agree except for those domains.
\end{definition}

\begin{remark}[what the full evaluation does and does not assert]
Definition~\ref{def:fulleval} is a total $\ZF$-definable function on syntax
and environments, for arbitrary ordinals $\lambda_i$: each clause forms
a set from sets by pairing, powerset, and the recursion theorem, and
$U_i$ is a set whatever $\lambda_i$ is. It asserts nothing: not that a
type lands in its universe, nor that a term lands in its type. Werner's
evaluation is definable in the same way for any $\kappa_i$; what needs
$\kappa_i$ inaccessible is the theorem that the evaluation is sound for
\emph{every} set environment, and \cite[Prop.~4.4]{paper1} shows that
this theorem is false below an inaccessible: a cofinal $f:X\to U_0$ is a
legitimate value for the context $X{:}\Type_0,\ f{:}X\to\Type_0$, and
$\sem{\Sigma x{:}X.\,f\,x}\rho$ is a code of rank $\lambda_0$, not in
$U_0$. Soundness is therefore a statement about a class of environments,
and \S\ref{sec:adm} is the choice of that class. The domain
$V_{\rk(c)\cdot\omega}$ of families has no semantic content; it makes a
family a total set function fixed by the code, so that codes exist
before the $\lambda_i$ and before admissibility are defined.
\end{remark}

\begin{definition}[hybrid majorants]\label{def:hyb}
Let $t$ be a term in context $\Gamma$, $\rho$ an environment assigning to
each variable an element of the full decoding of its type, and $\tau$ an
assignment to each big variable $F$ (Definition~\ref{def:smallbig}) of a
majorant $m_F\ge\hat\rho(F)$.
Define $M_t(\rho,\tau)$ by recursion on $t$:
\begin{itemize}
\item a big variable $F$ of the context $\Gamma$ has majorant $m_F$;
  every other variable, that is, a small variable of $\Gamma$ or a
  variable bound within $t$ whether small or big, has majorant
  $\widehat{\rho(x)}$, the canonical majorant of its actual value;
\item $M_{f\,a}:=M_f\bigl(\widehat{\sem a\rho}\bigr)$, the majorant of $f$
  at the canonical majorant of the \emph{actual} value of $a$;
\item $M_{\lambda x{:}A.b}:=(-,-,\varphi)$ with
  $\varphi(\chi):=\sup\{M_b(\rho[x{:=}a],\tau):a\in\El^{\mathrm f}(c),\ \hat a\le\chi\}$
  for $c:=\sem A\rho$, on the domain $\mathbb M^{(\rk c+1)}_{\rk(c)\cdot\omega}$
  (the domain that $\lhd_{\code\Pi(c,\cdot)}$ requires); every function
  part below has the domain that its syntactic type prescribes in the
  same way, and function majorants are applied by monotone extension;
\item the eliminators are functions, and their majorants are function
  parts that branch on the actual argument and are assembled from the
  majorants of the step terms. With $c$ the code of the argument type
  and $\chi$ ranging over the domain prescribed by $c$:
  \begin{itemize}
  \item $M_{\mathsf{ind}^{\mathbf n}_C\vec c}(\chi):=\sup\{M_{c_j}:j<n,\ j+1\le\chi_{\mathrm o}\}$;
  \item $M_{\mathsf{ind}^\N_C c\,f}(\chi):=\sup\{\mu_m:m<\omega,\ m+1\le\chi_{\mathrm o}\}$
    where $\mu_0:=M_c$ and $\mu_{m+1}:=M_f(\hat m)(\mu_m)$;
  \item $M_{\mathsf{ind}^\Sigma_C f}(\chi):=\sup\{M_f(\hat a)(\hat b):\langle a,b\rangle\in\El^{\mathrm f}(c),\ \widehat{\langle a,b\rangle}\le\chi\}$;
  \item $M_{\mathsf{ind}^W_C f}(\chi):=\sup\{\mu(t):t\in\El^{\mathrm f}(c),\ \hat t\le\chi\}$,
    where for a tree $t=\mathsf{tree}\,a\,d$ the majorant $\mu(t)$ is
    defined by recursion on the well-founded tree as
    $\mu(t):=M_f(\hat a)(\hat d)(\mu_d)$, with $\mu_d$ the function
    majorant $\chi'\mapsto\sup\{\mu(d(b)):\hat b\le\chi'\}$ on the domain
    prescribed by the code of the branching type at $a$;
  \item $M_{\mathsf{ind}^=_C f}(\chi_1)(\chi_2)(\chi_3):=\sup\{M_f(\hat a):a\in\El^{\mathrm f}(c),\ \hat a\le\chi_1,\ \hat a\le\chi_2\}$,
    the actual argument being a pair of equal elements and a proof;
  \item $M_{\mathsf{ind}^{\|\|}_C f}:=1$, the target being a proposition;
  \item $M_{\mathsf{ind}^/_C f\,h}(\chi):=\sup\{\inf_{a\in q}M_f(\hat a):q\in\El^{\mathrm f}(c),\ \hat q\le\chi\}$,
    the infimum being over the members of the class $q$, which are
    determined by the relation carried by the quotient code; it
    majorizes the common value $f(a)$ because $\lhd$ is closed under
    infima, while $M_f(\hat q)$ need not, the infimum $\hat q$ of the
    members' majorants not being attained in general;
  \item $M_{\mathsf{ind}^{\mathsf{Acc}}_C f}(\chi)(\chi'):=\sup\{\mu(x):x\in\El^{\mathrm f}(c)\text{ accessible for }\sem R\rho,\ \hat x\le\chi\}$,
    where $\mu(x):=M_f(\hat x)(\mu_x)$ by recursion along the accessible
    part of the set relation $\sem R\rho$, with $\mu_x$ the function
    majorant $\chi''\mapsto(-,-,1\mapsto\sup\{\mu(y):\sem R\rho(y,x)=\top,\ \hat y\le\chi''\})$
    on the domain prescribed by $c$, the inner function part taking
    the proof argument.
  \end{itemize}
  In every clause the ordinals occurring in the result are ordinals
  occurring in values of the step majorants $M_f$, $M_c$, $M_{c_j}$, by
  monotone extension; this is what the proofs of \S\ref{sec:closure}
  use;
\item a type former $\code Q(c,G)$ has majorant $\max(M_c,\ \sup_a M_{G}(\hat a))+3$,
  the rank bound of a pair of its components; quotients add one; propositional terms,
  $\mathsf{refl}$, $|{-}|$, proofs, and axioms have majorant $1$; and the
  majorant of a code as an element of a universe is its rank plus one.
\end{itemize}
\end{definition}

By induction on $t$, $\sem t\rho\lhd M_t(\rho,\tau)$ whenever every
bound variable met in the evaluation takes a value in the full decoding
of its type, in particular in the model (Lemma~\ref{lem:termadm}); in
the auxiliary structures of \S\ref{sec:closure} this can fail for the
recursive function of a $W$- or $\mathsf{Acc}$-recursion, and there the
majorants are only summed, not used as bounds (Lemma~\ref{lem:auxrank}).
The examples of
\S\ref{sec:examples} explain the two design choices, branching on actual
values and suprema over actual elements. One point deserves care: the
value of a term, and hence the branch an eliminator takes, can depend on
the truth values of universe-quantified propositions, in two distinct
ways. For a term-defined step of an $\mathsf{Acc}$-recursion this
dependence is only apparent: a recursive call at $y$ needs a proof of
$R\,y\,x$, which the step can obtain only from provable instances, true
in every model, or from a hypothesis $h$ in the context, true in every
environment that types the context. So for a fixed typed $\rho$ the
value is determined, and what varies between the $\alpha$-model and the
$\lambda$-model is whether $\rho$ is typed at all, since $\rho(h)$ has a
type in one and not the other. For an admissible element that is not the
denotation of a term the dependence is genuine: the step is an arbitrary
set function receiving the recursive-call function $\mathrm{rec}$ as an
argument, and $\mathrm{rec}(y)$ is the empty function or a value
function according to the truth of $R\,y\,x$, which a set function can
inspect. No term can do this, but the environments of
\S\ref{sec:adm} contain such elements. The majorant of the recursion
itself is unaffected in both cases; what is affected is a later
heterogeneous eliminator branching on the result. This is why the
closure functions below quantify over truth assignments and over
environments typed under each assignment, rather than over models.

\begin{definition}[truth assignments]\label{def:sigma}
A \emph{truth assignment} $\sigma$ assigns a truth value to every pair
$(\forall X{:}\Type_i.\,P,\ \rho)$ or $(\exists X{:}\Type_i.\,P,\ \rho)$ of
a universe-quantified proposition and an environment for its free
variables. The \emph{$\sigma$-evaluation} $\sem t^\sigma\rho$ is the
set-theoretic evaluation in which a universe-quantified proposition is
read off from $\sigma$ and every other construct is computed; the
\emph{actual} truth of the model is one particular $\sigma$. An
environment is a \emph{$\sigma$-environment for $\Gamma$} if each
$\rho(x)$ lies in $\El^{\mathrm f}(\sem{\Gamma(x)}^\sigma\rho)$. Hybrid
majorants $M^\sigma_t(\rho,\tau)$ are computed with $\sigma$-evaluation.
\end{definition}

An equation between codes is not an atom of $\sigma$: it is decided by
the codes. So a hypothesis $q:B(a)=A'$ in $\Gamma$ forces, under every
$\sigma$ for which $\rho$ is a $\sigma$-environment, the value
$\sem a^\sigma\rho$ to select the branch of $B$ that decodes to $A'$.
This is what lets the majorant branch on actual values without
depending on which $\sigma$ is the actual one.

\subsection{Auxiliary models and closure ordinals}\label{sec:closure}

\paragraph{Auxiliary models.}
This is the first phase: the extents are determined by recursion on a
level $N$, and at stage $N$ the ordinals $\lambda_0,\dots,\lambda_{N-1}$
are given and $\lambda_N$ is to be found.
For an ordinal $\alpha$ and an increasing sequence
$\vec\alpha'=(\alpha'_{N+1}<\alpha'_{N+2}<\cdots)$ of ordinals above $\alpha$, the
\emph{$(N,\alpha,\vec\alpha')$-model} has universes $U_j:=\{c:\rk c<\lambda_j\}$
for $j<N$, $U_N:=\{c:\rk c<\alpha\}$, and $U_j:=\{c:\rk c<\alpha'_j\}$ for
$j>N$. The lower universes are the true ones; the higher ones are
arbitrary, and the closure functions below take a supremum over
$\vec\alpha'$ as well, so that nothing at stage $N$ depends on
$\lambda_j$ for $j>N$. That this supremum is a set is
Lemma~\ref{lem:stab}. The auxiliary model is a structure over which
suprema are taken and is not claimed to be a model.

\begin{definition}[witness trees]\label{def:wtree}
An \emph{$N$-witness tree} $\theta$ for a term $t$ in context $\Gamma$
assigns to each $N$-big variable $F$ of $\Gamma$ that is free in $t$ a
term $t_F$ in some context $\Gamma_F$, having there a type of syntactic
value level at most $\vl(\Gamma(F))$, together with an $N$-witness tree
$\theta_F$ for $t_F$; the assignment is required to be well founded.
Being finitely branching and well founded, $\theta$ is a finite tree,
and there are countably many pairs $(t,\theta)$. The root $t$ may have
any type; the closure functions below are defined only when its value
level is at most $N$, while witness terms may have any value level, as
for $Y{:}\Type_{N+1}$. The bound on the value level of witnesses is
what makes the supremum over higher extents bounded in
Lemma~\ref{lem:stab}: with $\lambda Z.\,\Type_N$ admitted as a witness
for $F{:}\Type_{N+1}\to\Type_N$, the root $F\,Y$ would receive the
majorant $\alpha'_{N+1}+2$.
\end{definition}

\begin{definition}[witnessed environments]\label{def:wenv}
Fix $N$, an ordinal $\alpha$, higher extents $\vec\alpha'$, a truth
assignment $\sigma$ (Definition~\ref{def:sigma}), and a pair $(t,\theta)$
with $t$ in context $\Gamma$. A \emph{witnessed
$(N,\alpha,\vec\alpha',\sigma)$-environment for $(t,\theta)$} is a pair
$(\rho,\tau)$, defined by recursion on $\theta$, such that
\begin{itemize}
\item $\rho$ is a $\sigma$-environment for $\Gamma$ in the
  $(N,\alpha,\vec\alpha')$-structure: for every $x{:}A$ in $\Gamma$,
  $\rho(x)\in\El^{\mathrm f}(\sem A^\sigma\rho)$, where $\sem A^\sigma\rho$
  is a code by Definition~\ref{def:fulleval}, whether or not it lies in any
  universe;
\item $\tau$ assigns a majorant $\tau(F)$ to each $N$-big variable $F$ of
  $\Gamma$ free in $t$;
\item for each such $F$ there is a witnessed
  $(N,\alpha,\vec\alpha',\sigma)$-environment $(\rho_F,\tau_F)$ for
  $(t_F,\theta_F)$ with
  \[
  \tau(F)=M^\sigma_{t_F}(\rho_F,\tau_F)\qquad\text{and}\qquad
  \rho(F)\lhd_{\sem{\Gamma(F)}^\sigma\rho}\tau(F).
  \]
\end{itemize}
$N$-small variables, and $N$-big variables not free in $t$, are
constrained only by the first clause. The \emph{leaf supremum} of
$(\rho,\tau)$ is the leaf supremum of $M^\sigma_t(\rho,\tau)$; it is
well defined because $\tau$ is part of the data.
\end{definition}

Two facts are needed before a supremum over these environments can be
formed: that for fixed extents they form a set, and that the class of
truth assignments, which is a proper class of class functions, matters
only through a set.

\begin{lemma}[locality in $\sigma$]\label{lem:sigmaloc}
Fix $N,\alpha,\vec\alpha'$, a term $t$ in context $\Gamma$, and a set
$E$ of environments for $\Gamma$. There are sets $S_{t,E}$, of pairs of
a universe-quantified proposition and an environment, and $V_{t,E}$,
such that for every $\rho\in E$ and every truth assignment $\sigma$:
$\sem t^\sigma\rho\in V_{t,E}$; and if $\sigma'$ agrees with $\sigma$ on
$S_{t,E}$ then $\sem{t}^{\sigma'}\rho=\sem t^\sigma\rho$ and
$M^{\sigma'}_t(\rho,\tau)=M^\sigma_t(\rho,\tau)$ for every $\tau$.
\end{lemma}
\begin{proof}
By induction on $t$, defining $S_{t,E}$ and $V_{t,E}$ together.
\emph{Variables and constants:} $S:=\emptyset$; $V:=\{\rho(x):\rho\in E\}$,
or the constant's value. \emph{Universe-quantified propositions}
$Q=\forall X{:}\Type_i.\,P$ or $\exists X{:}\Type_i.\,P$: $S_{Q,E}:=\{(Q,\rho):\rho\in E\}$
and $V:=\{\bot,\top\}$; this is the only clause that consults $\sigma$.
\emph{Binders} $\Pi x{:}A.B$, and likewise $\Sigma$, $W$, quotients,
abstractions, and propositional quantifiers over non-universe types:
$V_{A,E}$ is a set by induction, so
$E':=\{\rho[x{:=}a]:\rho\in E,\ a\in V_{\rk(c)\cdot\omega},\ c\in V_{A,E}\}$
is a set; put $S:=S_{A,E}\cup S_{B,E'}$ and let $V$ consist of the
codes, functions, or truth values that the corresponding clause of
Definition~\ref{def:fulleval} forms from a member of $V_{A,E}$ and a
function into $V_{B,E'}$; this is a set. \emph{Application} $f\,a$:
$S:=S_{f,E}\cup S_{a,E}$ and $V:=\{F(A):F\in V_{f,E},\ A\in V_{a,E}\}\cup\{\emptyset\}$.
\emph{Pairs, constructors, and data:} componentwise. \emph{Eliminators:}
the recursion runs along an actual datum, which lies in the set
$V_{a,E}$ for the argument $a$, and the step functions are values in the
sets $V_{f,E}$ of the step terms; the environments the recursion
produces for the step terms therefore lie in a set $E''$ obtained from
$E$ by adjoining components, subtrees, predecessors, or members drawn
from members of $V_{a,E}$, and $S$ is the union of the $S_{f,E''}$ over
the step terms, $V$ the set of possible results. For $\mathsf{ind}^{\mathsf{Acc}}$
the relation $\sem R^\sigma$ varies with $\sigma$, but the recursion
visits only elements of the domain, a member of $V_{X,E}$, so $E''$ is a
set independent of $\sigma$. In each clause the evaluation consults
$\sigma$ only at pairs in $S$, so two assignments agreeing on $S$ give
the same value, by the same induction; and $M^\sigma_t(\rho,\tau)$ is
computed from $\sigma$-values of subterms at environments in these
sets, so it agrees as well.
\end{proof}

\begin{lemma}[witnessed environments form a set]\label{lem:envset}
Fix $N,\alpha,\vec\alpha'$ and $(t,\theta)$. There are a set $E_{t,\theta}$
of environments and a set $S_{t,\theta}$ of pairs such that every
witnessed $(N,\alpha,\vec\alpha',\sigma)$-environment $(\rho,\tau)$, for
any $\sigma$, has $\rho\in E_{t,\theta}$; for each $\rho\in E_{t,\theta}$
and each $\sigma$ the possible $\tau$ form a set; and
$M^\sigma_t(\rho,\tau)$ depends on $\sigma$ only through
$\sigma{\restriction}S_{t,\theta}$.
\end{lemma}
\begin{proof}
By recursion on $\theta$ and, inside it, on the length of $\Gamma$. For
the empty context $E:=\{\emptyset\}$. Given the set $E_{<x}$ of possible
restrictions of $\rho$ to the prefix before $x{:}A$, the codes
$\sem A^\sigma\rho$ for $\rho\in E_{<x}$ and all $\sigma$ lie in
$V_{A,E_{<x}}$ by Lemma~\ref{lem:sigmaloc}, so the possible values
$\rho(x)$ lie in $\bigcup\{\El^{\mathrm f}(c):c\in V_{A,E_{<x}}\}$, a
set, and $E_{\le x}$ is the set of extensions. For an $N$-big $F$ free in
$t$, the witnessed environments for $(t_F,\theta_F)$ have $\rho_F$ in a
set and, for each, $\tau_F$ in a set, by the recursion on $\theta$; so
the possible values $\tau(F)=M^\sigma_{t_F}(\rho_F,\tau_F)$ form a set.
Finally $S_{t,\theta}:=S_{t,E_{t,\theta}}\cup\bigcup_F S_{t_F,\theta_F}$.
\end{proof}

\begin{definition}[closure functions at level $N$]\label{def:nu}
Let $t$ have value level at most $N$ and let $\theta$ be an $N$-witness
tree for $t$. For an ordinal $\beta$ put
\[
\nu^{N,0}_{t,\theta}(\beta):=\sup\{\text{leaf supremum of }(\rho,\tau):\
\vec\alpha',\ \sigma,\ (\rho,\tau)\text{ a witnessed }(N,\beta,\vec\alpha',\sigma)\text{-environment for }(t,\theta)\},
\]
the supremum over all higher extents $\vec\alpha'$ above $\beta$, all
truth assignments $\sigma$, and all witnessed environments in the
$(N,\beta,\vec\alpha')$-structure, and
\[
\nu^N_{t,\theta}(\alpha):=\sup_{\beta\le\alpha}\nu^{N,0}_{t,\theta}(\beta).
\]
\end{definition}

\begin{lemma}[higher-input stability]\label{lem:stab}
Let $t$ have value level at most $N$. For fixed $\beta$ and $\vec\alpha'$
the leaf suprema of the witnessed $(N,\beta,\vec\alpha',\sigma)$-environments
for $(t,\theta)$, over all $\sigma$, are bounded by an ordinal
$\mu_{t,\theta}(\beta,\vec\alpha')$; and there is an ordinal
$B_{t,\theta}(\beta)$ with $\mu_{t,\theta}(\beta,\vec\alpha')\le B_{t,\theta}(\beta)$
for every $\vec\alpha'$.
\end{lemma}
The first half is Lemmas~\ref{lem:sigmaloc} and \ref{lem:envset}. The
restriction to value level at most $N$ is necessary for the second: the
term $\Type_{N+1}$ has majorant $\alpha'_{N+1}+2$ in the
$(N,\beta,\vec\alpha')$-structure. The proof of the second half is an
analogue, in the auxiliary structure, of the rank lemma of
Section~\ref{sec:rank}, and is prepared by a definition and a lemma.

\begin{definition}[content]\label{def:content}
Fix $N$ and an $(N,\beta,\vec\alpha')$-structure. For a code $c$ of
value level $\kappa(c)\le N$ and $e\in\El^{\mathrm f}(c)$, the
\emph{content} $\|e\|$ is the leaf supremum of the canonical majorant
$\hat e$; the content $\|c\|$ of the code is $\sup\{\|e\|:e\in\El^{\mathrm f}(c)\}$.
For a code $c$ with $\kappa(c)>N$ put $\|c\|:=0$ and $\|e\|:=0$ for its
elements. An environment $\rho$ for a context $\Gamma$ has content
$<\gamma$ if $\|\rho(x)\|<\gamma$ for every variable $x$ whose type has
syntactic value level at most $N$.
\end{definition}
Unwinding Definition~\ref{def:codesbd}: for a code as an element of a
universe the content is its rank plus one; for a function it is the
supremum of the contents of its values, the inputs contributing
nothing; for a pair or a class the maximum, respectively infimum, over
components or members; for a tree the maximum of its height plus one
and the contents of its labels. Since $\El^{\mathrm f}(c)\subseteq V_{\rk(c)\cdot\omega}$,
$\|c\|\le\rk(c)\cdot\omega$; and $\|c\|<\beta\cdot\omega$ for $c$ in any
universe $U_j$, $j\le N$, of the structure, while for $\code{\U_j}$
itself $\|\code{\U_j}\|\le\beta+1$ if $j=N$ and $\le\lambda_j+1$ if
$j<N$.

\begin{lemma}[rank lemma in the auxiliary structure]\label{lem:auxrank}
Fix $N$ and $\beta$. For every term $t$ in context $\Gamma$ with
$\vl(t)\le N$ there is a $\ZF$-definable ordinal function
$g_t:\mathrm{On}\to\mathrm{On}$, depending only on $t$, $\Gamma$, $\beta$,
and $\lambda_0,\dots,\lambda_{N-1}$, such that for every
$\vec\alpha'$, every $\sigma$, and every $\sigma$-environment $\rho$ for
$\Gamma$ in the $(N,\beta,\vec\alpha')$-structure of content $<\gamma$:
$\|\sem t^\sigma\rho\|<g_t(\gamma)$, the content being taken in the
code $\sem T^\sigma\rho$ of $t$'s type when $\sem t^\sigma\rho$ lies in
its decoding, and the statement being vacuous otherwise.
\end{lemma}
\begin{proof}
By induction on $t$, with $g_t$ monotone; all bounds below are uniform
in $\vec\alpha'$ and $\sigma$ because nothing in them mentions a
universe above $N$ except as an index set.
\emph{Variables:} $\|\rho(x)\|<\gamma$ by hypothesis; $g_x(\gamma):=\gamma$.
\emph{Constants:} finite, or $\lambda_j+2$ for $\Type_j$ with $j<N$
(the term $\Type_N$ has value level $N+1$).
\emph{Application $f\,a$}, with $f:\Pi x{:}A.B$ and $\vl(B)\le N$: the
value is $\sem f\rho(\sem a\rho)$, whose content is at most the content
of the function $\sem f\rho$, which is the supremum of the contents of
its values; so $g_{f\,a}:=g_f$. (The argument's content, which may be
undefined or huge when $\vl(A)>N$, is not used.)
\emph{Abstraction $\lambda x{:}A.\,b$:} the content of the function is
$\sup_a\|\sem b\rho[x{:=}a]\|$ over $a\in\El^{\mathrm f}(\sem A\rho)$. If
$\vl(A)>N$ the extended environment has the same content and
$g:=g_b$. If $\vl(A)\le N$ then $\|a\|\le\|\sem A\rho\|<g_A(\gamma)$ by
induction ($A$ is a type term of value level $\le N$, hence a root), so
$g_{\lambda x.b}(\gamma):=g_b(\max(\gamma,g_A(\gamma)))$.
\emph{Pairs, constructors, numerals, classes:} the maximum of the
components' bounds, or $\omega$.
\emph{Type formers} $\Pi x{:}A.B$, $\Sigma$, $W$, $A/R$ at sort
$\Type_i$, $i\le N$: the value is the code $\code Q(\sem A\rho,G)$ with
$G(a)=\sem B\rho[x{:=}a]$ or $\bot$; its content is its rank plus one,
and its rank is at most $\max(\rk\sem A\rho,\sup_a\rk G(a))+3$. Both
$\rk\sem A\rho$ and $\rk G(a)$ are contents, of $\sem A\rho$ as an
element of $\code{\U_j}$ and of $G(a)$ as an element of $\code{\U_k}$,
bounded by $g_A(\gamma)$ and $g_B(\max(\gamma,g_A(\gamma)))$
respectively, the latter uniformly in $a$ since $\|a\|<g_A(\gamma)$
when $\vl(A)\le N$ and $a$ is irrelevant to the content otherwise.
\emph{Propositions, proofs, axioms:} content $1$.
\emph{Eliminators for $\mathbf n$, $\N$, $\Sigma$, $=$, $\|\cdot\|$, and
quotients:} the value is a step function applied to components of the
datum and, for $\N$, to the value at the predecessor; contents are
bounded by $g$ of the step term iterated finitely often, or by the
component's content, which is at most the datum's content (or
irrelevant if the datum has value level $>N$, in which case only its
level-$\le N$ components are used, and those have content bounded by
the type terms of the components, by induction). For $\mathsf{ind}^\N$
the bound is $\sup_n g_s^{n}(\gamma)$ over the step $s$, an ordinal.
\emph{$W$-recursion} $\mathsf{ind}^W s\,w$ with $w:Wx{:}A.B$ and step
$s=\lambda x\,d\,r.\,b$ (or a variable, which is the trivial case): the
value at a tree $w$ is $\sem b$ at the environment extended by the label,
the subtree function, and the recursive function $r_w$, whose values
are the values at the subtrees. By induction on the tree, if the
values at all subtrees have content $<\delta$ then the value at $w$ has
content $<g_b(\max(\gamma,\delta,\ \|\text{label}\|,\ \|d\|))$, where
the label's content is $<g_A(\gamma)$ if $\vl(A)\le N$ and irrelevant
otherwise, and $\|d\|$ is the supremum of the subtrees' contents, which
is at most the tree's content. Let $h(w)$ be the height of $w$. Then
the content of the value at $w$ is below the $h(w)$-th iterate of the
monotone function $\delta\mapsto g_b(\max(\gamma,\delta,g_A(\gamma),\|w\|))$,
and $\|w\|$ is bounded by $g$ of the type term $Wx{:}A.B$, which is a
root of value level at most $N$ when $\vl(W x{:}A.B)\le N$; when the
tree's type has value level $>N$ (a big label type $A$), its height is
still bounded: every node has children indexed by $\El^{\mathrm f}(\sem B\rho[x{:=}a])$,
a set of rank below $g_B(\gamma')\cdot\omega$ for
$\gamma':=\max(\gamma,g_A(\gamma))$ by induction, so
$h(w)<\hartogs(V_{g_B(\gamma')\cdot\omega})$, the labels contributing
nothing to the height. So $g_{\mathsf{ind}^W s\,w}$ is the corresponding
iterate, an ordinal depending on $\gamma$ only.
\emph{$\mathsf{Acc}$-recursion} $\mathsf{ind}^{\mathsf{Acc}}s\,x\,p$ along
$R$ on $X$ with step $s=\lambda x\,r.\,b$: the value at $x$ is $\sem b$
at the environment extended by $x$ and the recursive function $r_x$ on
the predecessors of $x$; $\sem b$ applies $r_x$ only at the values of
the finitely many subterms of $b$ of type $X$ in the finitely many
environments the evaluation of $b$ produces, and those environments
differ from the given one only in bound variables of $b$ whose types
have value level $\le N$ (a bound variable of higher value level would
make the containing subterm of value level $>N$, and such subterms are
not evaluated for their value but only for their truth value or as
inputs). So the set of predecessors at which $r_x$ is inspected has
size at most that of $V_{\delta}$ for $\delta:=\sup$ of the $g$-bounds
of those bound variables' types, and the tree of inspected predecessors
below $x$, which is well founded because $x$ is accessible, has height
below $\hartogs(V_{\delta+\omega})$. The value at $x$ depends only on
$r_x$ at inspected predecessors, so by induction along that tree its
content is below the $\hartogs(V_{\delta+\omega})$-th iterate of
$\delta'\mapsto g_b(\max(\gamma,\delta'))$, where the content of $x$
itself is irrelevant if $\vl(X)>N$ and $<g_X(\gamma)$ otherwise. When
$s$ is a variable its values lie in the cut decodings and the trivial
bound applies.
\end{proof}

The two recursion cases are where the auxiliary structure differs from
a model: the recursive function $r_w$ or $r_x$ has values outside the
cut universe when the recursion climbs, so it is not an element of the
cut decoding of its type, and the majorant $M^\sigma_t$, which sums over
elements of the cut decodings, does not majorize the value. The lemma
bounds the value anyway, by the height of the recursion, exactly as
$\wfrec$ is bounded in Lemma~\ref{lem:rank}. This unsoundness of the
majorants in the auxiliary structure is harmless: the first phase only
takes suprema of majorants, and in the second phase the transfer is
applied only at cuts above every reached rank.

\begin{proof}[Proof of the second half of Lemma~\ref{lem:stab}]
We prove, by induction on the number of nodes of $\theta$ and, for a
fixed number, on $|\Gamma|+|t|$, the following stronger statement:
there is an ordinal $B_{t,\theta}(\beta)$ such that for all
$\vec\alpha'$, $\sigma$, all witnessed $(N,\beta,\vec\alpha',\sigma)$-environments
$(\rho,\tau)$ for $(t,\theta)$ on a prefix of $\Gamma$, and all
$\sigma$-extensions $\rho'$ of $\rho$ to the rest of $\Gamma$ by
arbitrary values, with unwitnessed variables carrying the canonical
majorants of their values, every leaf of $M^\sigma_t(\rho',\tau)$ is
below $B_{t,\theta}(\beta)$. This is the form in which the induction
passes under binders. Write $\gamma$ for a bound on the content of
$\rho'$: every value $\rho'(x)$ whose type $A$ has $\vl(A)\le N$ lies in
$\El^{\mathrm f}(\sem A\rho')$, and $A$ is a type term of value level
$\le N$ in a shorter context, so $\|\rho'(x)\|<g_A(\gamma_{<x})$ by
Lemma~\ref{lem:auxrank} applied to the prefix; by induction on the
context, all contents are below an ordinal $\gamma_\Gamma(\beta)$
depending only on $\Gamma$ and $\beta$. Now the cases.
\emph{Small or unwitnessed variables:} the majorant is $\widehat{\rho'(x)}$,
whose leaves are below $\|\rho'(x)\|<\gamma_\Gamma$.
\emph{Witnessed variables $F$ as the root:} $\vl(F)\le N$, the majorant
is $M^\sigma_{t_F}(\rho_F,\tau_F)$, and the induction hypothesis for
$(t_F,\theta_F)$, with fewer nodes, applies.
\emph{Constants:} finite or $\lambda_j+2$, $j<N$.
\emph{Application $f\,a$:} $M_{f\,a}=M_f(\widehat{\sem a\rho'})$ by
monotone extension, whose leaves are among the leaves of $M_f$; and
$\vl(f)=\vl(B)\le N$ with $|f|<|t|$. The argument contributes only
through its canonical majorant as an input, never through its leaves.
\emph{Abstraction $\lambda x{:}A.\,b$:} the function part is
$\chi\mapsto\sup\{M^\sigma_b(\rho'[x{:=}a],\tau):a\in\El^{\mathrm f}(\sem A\rho'),\ \hat a\le\chi\}$;
its leaves are leaves of the majorants $M^\sigma_b(\rho'[x{:=}a],\tau)$,
which are instances of the induction hypothesis for $b$ in the context
$\Gamma,x{:}A$ (the bound variable being an unwitnessed extension), with
$|\Gamma,x{:}A|+|b|<|\Gamma|+|\lambda x{:}A.\,b|$; the bound $B_{b,\theta}(\beta)$
is uniform in $a$. Inputs $\chi$ are not leaves.
\emph{Pairs, constructors, and data:} componentwise; a component of a
term of value level $\le N$ has value level $\le N$, by the clauses for
$\vl$.
\emph{Type formers at sort $\Type_i$, $i\le N$:} the majorant is
$\max(M_A,\sup_aM_{\lambda x.B}(\hat a))+3$, whose leaves are among those
of $M_A$ and $M_{\lambda x.B}$, both terms of value level $\le i$ and
smaller than $t$.
\emph{Propositions, proofs, axioms:} $1$.
\emph{Eliminators:} in every clause of Definition~\ref{def:hyb} the
majorant of an eliminator is assembled from the majorants of the step
terms, applied to canonical majorants of actual data and to
majorants of recursive calls, together with infima and suprema over
actual members and branches; by monotone extension all its leaves are
leaves of the step terms' majorants, and the step terms have value
level $\le N$ (their type ends in the motive) and are smaller than
$t$. The recursion along a tree or an accessible part does not
accumulate: each node's majorant is a value of the step majorant, and
the induction hypothesis bounds all values of the step majorant at
once. (This is where the hybrid definition matters: a syntactic
majorant that added a constant at every node would grow with the
height of the datum, which is unbounded in $\vec\alpha'$ for a datum of
higher level.)
This exhausts the clauses, and $B_{t,\theta}(\beta)$ is the maximum of
the bounds used, which depend only on $t$, $\Gamma$, $\theta$, and
$\beta$.
\end{proof}

The case that the earlier sketch of this lemma missed is the recursion
case of Lemma~\ref{lem:auxrank}: a term of value level $N$ can use
higher-level data directly, through a recursor, and not only through
propositions and applications, and the \emph{value} of such a
recursion is not bounded by its majorant in the auxiliary structure.
What bounds the value is the height of the part of the datum that a
term can inspect, which is indexed by level-$\le N$ types; and what
bounds the majorant is that it is assembled from step majorants whose
leaves are bounded uniformly over all inputs.

\begin{lemma}\label{lem:nuset}
By Lemma~\ref{lem:stab}, for every $(t,\theta)$ with $t$ of value
level at most $N$: $\nu^{N,0}_{t,\theta}(\beta)$ and
$\nu^N_{t,\theta}(\alpha)$ are ordinals; $\nu^N_{t,\theta}$ is monotone
and $\nu^N_{t,\theta}\ge\nu^{N,0}_{t,\theta}$; and $\alpha\mapsto\nu^N_{t,\theta}(\alpha)$
is $\ZF$-definable from $\lambda_0,\dots,\lambda_{N-1}$, without
reference to $\lambda_j$ for $j\ge N$.
\end{lemma}
\begin{proof}
For fixed $\beta$ and $\vec\alpha'$ the leaf suprema range over a set:
$\rho$ over $E_{t,\theta}$, $\tau$ over a set for each $\rho$, and
$\sigma$ through the set $2^{S_{t,\theta}}$ of restrictions, by
Lemma~\ref{lem:envset}. So their supremum $\mu_{t,\theta}(\beta,\vec\alpha')$
is an ordinal, and by Lemma~\ref{lem:stab} the class
$\{\mu_{t,\theta}(\beta,\vec\alpha'):\vec\alpha'\}$ is bounded, so
$\nu^{N,0}_{t,\theta}(\beta)$ is an ordinal; then $\nu^N_{t,\theta}(\alpha)$
is a supremum of $\alpha+1$ ordinals. Monotonicity and the inequality
hold by construction. For definability, every quantifier in
Definitions~\ref{def:wenv} and \ref{def:nu} ranges over a set
constructed by the recursions of Lemmas~\ref{lem:sigmaloc} and
\ref{lem:envset}, except the one over $\vec\alpha'$, which is a
supremum of a class function of $\vec\alpha'$ shown bounded by
Lemma~\ref{lem:stab} and hence a set supremum; the parameters are the
lower universes and their universe constants.
\end{proof}

\begin{remark}[why the supremum over $\beta\le\alpha$]
The inner function $\nu^{N,0}_{t,\theta}$ need not be monotone, and the
naive argument ``more codes give more environments'' is wrong. An
$N$-big value $F$ on the domain $U_N^\beta$ is a set function on the
codes of rank below $\beta$; in the $(N,\alpha,\vec\alpha')$-structure
for $\alpha>\beta$ its type decodes to functions on the codes of rank
below $\alpha$, of which $F$ is not one. So the witnessed environments
at $\beta$ are not among those at $\alpha$. For the same reason
$\nu^{N,0}$ has no continuity at limits: a value on $U_N^{\alpha_\omega}$
has no counterpart at any $\alpha_n$, short of the retraction of
\S\ref{sec:adm}, which the first phase does not have. The closure
argument below uses monotonicity only, so it is built into
$\nu^N_{t,\theta}$ by definition, at no cost: the transfer lemma
delivers a bound of the form $\nu^{N,0}_{t,\theta}(\alpha)$, and
$\nu^{N,0}\le\nu^N$.
\end{remark}

Only $N$-big variables are restricted by witnesses, because the transfer
at level $N$ (Lemma~\ref{lem:transfer}) fixes everything of level below
$N$.

\begin{definition}[closed ordinals]\label{def:closed}
Given $\lambda_0<\dots<\lambda_{N-1}$, an ordinal $\lambda>\lambda_{N-1}$
(or $\lambda>\omega$ if $N=0$) is \emph{$N$-closed} if it is a limit,
$\hartogs(V_{\beta+1})<\lambda$ for every $\beta<\lambda$, and
$\nu^N_{t,\theta}(\beta)<\lambda$ for every $\beta<\lambda$ and every pair
$(t,\theta)$ with $t$ of value level at most $N$. Then $\lambda_N$ is the
least $N$-closed ordinal.
\end{definition}

\begin{lemma}\label{lem:closedex}
An $N$-closed ordinal exists, so
$\lambda_N$ is defined; and $\lambda_N$ has cofinality $\omega$.
\end{lemma}
\begin{proof}
Enumerate the countably many pairs as $(t_k,\theta_k)_{k\in\omega}$.
Put $\alpha_0:=\lambda_{N-1}+1$ (or $\omega+1$) and
\[
\alpha_{n+1}:=\sup\bigl(\{\hartogs(V_{\alpha_n+1})\}\cup\{\nu^N_{t_k,\theta_k}(\alpha_n):k\le n\}\bigr)+1,
\]
a successor of a finite supremum of ordinals by Lemma~\ref{lem:nuset},
and $\lambda:=\sup_n\alpha_n$. The sequence is strictly increasing, so
$\lambda$ is a limit of cofinality $\omega$ above $\lambda_{N-1}$. Let
$\beta<\lambda$ and $k\in\omega$; choose $n\ge k$ with $\beta<\alpha_n$.
Then $\nu^N_{t_k,\theta_k}(\beta)\le\nu^N_{t_k,\theta_k}(\alpha_n)<\alpha_{n+1}<\lambda$
by monotonicity (Lemma~\ref{lem:nuset}), and
$\hartogs(V_{\beta+1})\le\hartogs(V_{\alpha_n+1})<\alpha_{n+1}<\lambda$.
So $\lambda$ is $N$-closed, and the least $N$-closed ordinal exists.
For its cofinality, run the same iteration from
$\alpha_0:=\lambda_{N-1}+1$ but note that, $\lambda_N$ being $N$-closed
and a limit, each $\alpha_n<\lambda_N$ by induction; so the resulting
closed ordinal $\lambda'=\sup_n\alpha_n$ satisfies $\lambda'\le\lambda_N$,
hence $\lambda'=\lambda_N$ by leastness, and $\lambda_N$ has cofinality
$\omega$.
\end{proof}

So the sequence $(\lambda_N)$ is defined by recursion on $N$, each
stage using only the full decoding, the majorants, and the previous
stages; nothing about admissibility or the retraction is presupposed,
and nothing about $\lambda_j$ for $j>N$, by Lemma~\ref{lem:nuset}. In
$\ZFC$ the $\lambda_N$ are strong limit cardinals by
Example~\ref{ex:set}. The proofs above use of the eliminator clauses of
Definition~\ref{def:hyb} only that every leaf of an eliminator's
majorant is a leaf of a step majorant, which the clauses make explicit.

\begin{example}[the closure functions saturate, and the closure condition is too weak]\label{ex:sat}
The lemmas of this subsection are correct, but the ordinals they
produce are not the right ones. In the structure with $U_N$ cut at
$\alpha$, a bound variable of universe type ranges over codes of rank
below $\alpha$, and the recursive function of an $\mathsf{Acc}$- or
$W$-recursion, which is a bound variable of the step term, ranges over
functions into the cut universe. So the majorant of a rank-raising
recursion saturates at about $\alpha+3$ whatever the rank of its actual
value: this is the unsoundness of the majorants in the auxiliary
structure noted after Lemma~\ref{lem:auxrank}, and it is not harmless.
It makes $\nu^N_{t,\theta}(\alpha)$ about $\alpha$ plus a small ordinal
for every $(t,\theta)$, so closure under the $\nu^N_{t,\theta}$ is nearly
vacuous and $\lambda_N$ is the least ordinal $\eta>\lambda_{N-1}$ closed
under $\alpha\mapsto\hartogs(V_{\alpha+1})$.

For that $\lambda_N$ the type-formation lemma is false. For
$X:\Type_N$ let $\mathrm{WO}(X)$ be the type of well-founded linear
orders $R$ on $X$, and for $o=(R,\dots):\mathrm{WO}(X)$ define
$T_o:X\to\Type_N$ by $\mathsf{Acc}$-recursion with motive $\Type_N$,
\[
T_o(x):=\bigl(\Sigma y{:}X.\,\Sigma h{:}R\,y\,x.\,T_o(y)\bigr)\to\Prop ,
\]
which is well typed at $\Type_N$, and put
$t(X):=\Sigma o{:}\mathrm{WO}(X).\,\Sigma x{:}X.\,T_o(x):\Type_N$. Then
$\rk T_o(x)$ is at least the order type of $x$ in $R$, so
$\rk t(X)\ge\hartogs(\El(X))$; and $\El(T_o(x))$ is a cumulative
hierarchy of height the order type of $x$, so the ranks $r_n$ of
$t^n(\N)$ satisfy $r_{n+1}\ge\hartogs(V_{r_n})$. Each $t^n(\N)$ has rank
below $\eta$ and lies in $U_N$, and $n\mapsto t^n(\N)$ is a term, by
$\mathsf{ind}^\N$ with motive $\Type_N$. But
$\Sigma n{:}\N.\,t^n(\N):\Type_N$ has rank $\sup_nr_n=\eta=\lambda_N$ and
is not in $U_N$. Lemma~\ref{lem:transfer} does not apply to it, since no
cut below $\lambda_N$ lies above the ranks reached; and no repair of
that lemma can help, because the defect is in the definition of
$\lambda_N$.

In Section~\ref{sec:rank} the same iteration is harmless because the
majorant $f_t$ is a total ordinal function, computed in the universe
and not in a cut of it, so that closure under the majorant of
$n\mapsto t^n(\omega)$ puts $\sup_nf_t^n(\omega)$ below $\lambda$. The
closure functions here must likewise not saturate.
\end{example}

\begin{remark}[the repair, not carried out]\label{rem:nosat}
Take the supremum over the cut as well. For an ordinal $\gamma$ let
$\bar\nu^N_{t,\theta}(\gamma)$ be the supremum, over \emph{all} cuts
$\alpha\ge\gamma$, all $\vec\alpha'$, all $\sigma$, and all witnessed
$(N,\alpha,\vec\alpha',\sigma)$-environments for $(t,\theta)$, of the
$\gamma$-leaves $\ell_\gamma(M^\sigma_t(\rho,\tau))$ of
Definition~\ref{def:relmaj}, made monotone in $\gamma$ as before, and
let $\lambda_N$ be the least limit ordinal above $\lambda_{N-1}$ closed
under Hartogs numbers and under every $\bar\nu^N_{t,\theta}$. The
$\gamma$-leaves are needed because the full leaf supremum of
$M_{\lambda X{:}\Type_N.X}$ is the cut itself. In a cut above every rank
reached the majorants are sound and do not saturate, so
$\bar\nu^N$ of the term $\Sigma n{:}\N.\,t^n(\N)$ is at least $\eta$ and
$\lambda_N>\eta$, as it must be; and no term can diagonalize against
$\lambda_N$, whose defining sequence uses the closure functions of all
terms at once.

Two consequences. First, the definition now needs a \emph{cut-stability
lemma}: for fixed $(t,\theta)$ and $\gamma$ the $\gamma$-leaves are
bounded as the cut grows. We expect this to hold for the reason
Lemma~\ref{lem:stab} holds in the higher extents: witnessed environments
are small whatever the cut, a bound variable summed under
$\hat a\le\chi$ with $\chi$ hereditarily below $\gamma$ takes values
below $\gamma$, and a term applies its function arguments only at
values it has reached, so nothing at small inputs depends on the
cut except through truth values, over which the supremum is already
taken. The proof should be a cut-free form of Lemma~\ref{lem:auxrank},
in which the bound on a function majorant is a function of the bounds
on its inputs; it has not been written. Second, the model's own
structure, with $U_N$ cut at $\lambda_N$ and the true higher universes,
is then \emph{one of the structures over which $\bar\nu^N$ takes its
supremum}, and an admissible environment with its witnesses is a
witnessed environment of it, since
$M^{\mathcal A}_{t_F}\le M_{t_F}$ and $\lhd$ is upward closed. So for a
family $\lambda x.B$ over a code $c\in U_N$ and
$\gamma>\rk(c)\cdot\omega$,
\[
\sup_{a\in\El(c)}\rk\sem B\rho[x{:=}a]\ \le\ \ell_\gamma\bigl(M_{\lambda x.B}(\rho,\tau)\bigr)\ \le\ \bar\nu^N_{\lambda x.B,\theta}(\gamma)\ <\ \lambda_N
\]
directly, by soundness of the majorants in the model and closure.
No transfer between structures is needed, hence no retraction, no
bounded support, and no comparison of majorants across cuts: all of
\S\ref{sec:transfer}, Definition~\ref{def:retr}, the support clause of
Definition~\ref{def:adm}, Lemma~\ref{lem:bspres}, and
Remark~\ref{rem:gap} would be deleted, admissibility would mean
hereditarily majorized by a term majorant and nothing else, and the
element $H$ of Example~\ref{ex:probe} would simply be admissible, which
does no harm once nothing has to commute with a retraction. This is
the shape of the argument of Section~\ref{sec:model}: a rank lemma
proved in the universe, and a model cut out at a closure point. The
whole weight of the construction would then rest on the cut-stability
lemma.
\end{remark}

\subsection{The retraction, admissible elements, and the model}\label{sec:adm}

\begin{definition}[retraction at level $N$]\label{def:retr}
From here on all the extents $\lambda_j$ are fixed by \S\ref{sec:closure};
the retraction plays no part in determining them. For a level $N$ and
an ordinal $\alpha$ define by recursion on codes maps $r^N_\alpha$ from
codes and elements of the model to those of the
$(N,\alpha,\vec\lambda_{>N})$-model, whose higher universes are the true
ones, and $e^N_\alpha$ in the other direction:
\begin{itemize}
\item both are the identity on constants and on every code and element
  of level other than $N$ that has no component at level $N$; on a
  universe element $X\in U_N$, $r^N_\alpha X:=X$ if $\rk X<\alpha$ and
  $\bot$ otherwise, and $e^N_\alpha$ is the inclusion;
\item $r^N_\alpha\code\Pi(c,G):=\code\Pi(r^N_\alpha c,\ a'\mapsto r^N_\alpha G(e^N_\alpha a'))$
  and likewise for $\Sigma$ and $W$; for quotients
  $r^N_\alpha\code/(c,R):=\code/(r^N_\alpha c,\ R')$ with
  $R':=\{(r^N_\alpha x,r^N_\alpha y):(x,y)\in R\}$, the image of $R$;
  $e^N_\alpha$ on codes extends a family by $\bot$ off the image of
  $e^N_\alpha$ and pulls a relation back, $R'':=\{(x,y):(r^N_\alpha x,r^N_\alpha y)\in R'\}$;
\item $r^N_\alpha(F)(a'):=r^N_\alpha\bigl(F(e^N_\alpha a')\bigr)$ for a function
  $F$ on $\El^{\mathrm f}(c)$ and $a'\in\El^{\mathrm f}(r^N_\alpha c)$ in the
  $(N,\alpha)$-model; $e^N_\alpha(F')(a):=e^N_\alpha\bigl(F'(r^N_\alpha a)\bigr)$;
  pairs and trees componentwise, $r^N_\alpha\langle a,b\rangle:=\langle r^N_\alpha a,r^N_\alpha b\rangle$
  and $r^N_\alpha(\mathsf{tree}\,a\,d):=\mathsf{tree}\,(r^N_\alpha a)\,(b'\mapsto r^N_\alpha d(e^N_\alpha b'))$,
  with $e^N_\alpha$ likewise; classes by $r^N_\alpha[x]_R:=[r^N_\alpha x]_{R'}$
  and $e^N_\alpha[x']_{R'}:=[e^N_\alpha x']_R$; truth values unchanged.
\end{itemize}
These clauses define $r^N_\alpha$ and $e^N_\alpha$ as set maps by
recursion on codes, with no side condition, once $e^N_\alpha$ on a class
is taken to be $\emptyset$ when $[e^N_\alpha x']_R$ depends on the
representative $x'$. So there is no circularity with
Definition~\ref{def:adm}, where bounded support is defined from these
maps. What does need a condition is that $r^N_\alpha$ maps
$\El^{\mathrm f}(c)$ into $\El^{\mathrm f}(r^N_\alpha c)$ and that
$r^N_\alpha e^N_\alpha=\mathrm{id}$; this is proved by induction on codes
for the admissible codes and elements of \S\ref{sec:adm}, the step at a
code using bounded support only of its proper constituents. For a pair
$\langle a,b\rangle$ with $b\in\El^{\mathrm f}(G(a))$, the retracted
second component must lie in $\El^{\mathrm f}(r^N_\alpha G(e^N_\alpha r^N_\alpha a))$,
the fibre of the retracted code at $r^N_\alpha a$, and $e^N_\alpha r^N_\alpha a$
need not be $a$; the two fibres agree when $r^N_\alpha G(a)$ depends on
$a$ only through $r^N_\alpha a$, which is the bounded support of the
family $G$ at $\alpha$ above its support. The same applies to the
subtree function of a tree. For a class, $r^N_\alpha$ is well defined
by construction of $R'$, and $e^N_\alpha$ is well defined when
$x\sim_R e^N_\alpha r^N_\alpha x$ for all $x$; this holds when the
relation $R$, as a function into $\Prop$, has bounded support at
$\alpha$ above its support, since then $R(x,y)=R(e^N_\alpha r^N_\alpha x,y)$
for all $y$, and $\sim_R$ is reflexive. On codes of rank below
$\alpha$ and their elements, $r^N_\alpha$ and $e^N_\alpha$ are the
identity, by induction on codes, and no condition is needed.
So $r^N_\alpha$ acts on a code or element of level above $N$ only by
restricting, through its families and arguments, everything at level
$N$ to rank below $\alpha$; it never collapses a higher-level object.
Then $r^N_\alpha e^N_\alpha=\mathrm{id}$; $r^N_\alpha$ is the identity on
elements of hereditary rank below $\alpha$ and on elements of level
below $N$; $\widehat{r^N_\alpha a}\le\hat a$; and the retractions are
compatible, $r^N_\alpha=r^N_\alpha\circ r^{N'}_{\alpha'}$ for $N\le N'$,
which one proves by a simultaneous induction on codes for $r$ and $e$,
using that both extend by $\bot$ uniformly, so that
$e^{N'}e^{N}=e^{N}$ on values.
\end{definition}

\begin{example}[an element with unbounded support]\label{ex:probe}
Let $H:(\Type_0\to\Type_0)\to\N$ be $H(G):=0$ if $G$ is the identity on
all of $U_0$ and $1$ otherwise. $H$ is majorized by the constant $2$, the
majorant of $\lambda G.\,1$, so without a support condition it would be
admissible. But for every $\alpha<\lambda_0$ and every extension
$e^i_\alpha$ of $\alpha$-elements to $\lambda$-elements,
$e^i_\alpha r^i_\alpha(\mathrm{id})\ne\mathrm{id}$, so $H(e^i_\alpha r^i_\alpha(\mathrm{id}))=1\ne H(\mathrm{id})$
and no retraction commutes with the application $H(\mathrm{id})$. A
term cannot behave like $H$: $\sem{\lambda G.\,b}\rho$ applies $G$ only
to arguments that $b$ builds, of rank bounded by the majorants. This is
the bounded-support condition of Definition~\ref{def:adm}, and the
seventh correction.
\end{example}

\begin{definition}[relative majorants and $\alpha$-leaves]\label{def:relmaj}
For a class $\mathcal X$ of sets, the \emph{$\mathcal X$-majorant}
$M^{\mathcal X}_t(\rho,\tau)$ is defined by the clauses of
Definition~\ref{def:hyb} with one change: the abstraction clause sums
only over arguments in $\mathcal X$,
\[
M^{\mathcal X}_{\lambda x.b}(\rho,\tau)(\chi):=\sup\{M^{\mathcal X}_b(\rho[x{:=}a],\tau):a\in\El^{\mathrm f}(\sem A\rho)\cap\mathcal X,\ \hat a\le\chi\}.
\]
So $M^{\mathcal X}_t\le M_t$ pointwise, $M^{\mathcal X}_t$ is monotone in
$\mathcal X$, and $M^V_t=M_t$. For an ordinal $\alpha$ the
\emph{$\alpha$-leaves} of a majorant $m$ are: its ordinal part; the
$\alpha$-leaves of the components of its pair part; and the
$\alpha$-leaves of $m_{\mathrm f}(\chi)$ for those $\chi$ in the domain
of the function part that are hereditarily below $\alpha$. Their
supremum is $\ell_\alpha(m)$, and $\ell_\alpha(m)\le\ell(m)$, the leaf
supremum.
\end{definition}

The $\alpha$-leaves are what the type-formation lemma needs, and the
distinction is not cosmetic. In the model, $M_{\lambda X{:}\Type_i.X}$ has
the value $\sup\{\rk X+1:X\in U_i,\ \hat X\le\chi\}=\lambda_i$ at every
input $\chi\ge\lambda_i$ in its domain, so its leaf supremum is
$\lambda_i$; in the structure with $U_i$ cut at $\alpha$ the same
majorant has leaf supremum $\alpha$. No retraction can make the two
leaf suprema equal, and an earlier version of Lemma~\ref{lem:transfer}
claimed exactly that. What is true is that the two agree at inputs
hereditarily below $\alpha$, and a type family $\lambda x.B$ over a
code $c\in U_i$ is only ever applied to inputs $\hat a$ with
$a\in\El^{\mathrm f}(c)\subseteq V_{\rk(c)\cdot\omega}$, hereditarily
below any $\alpha>\rk(c)\cdot\omega$.

\begin{definition}[admissibility]\label{def:adm}
An environment $\rho$ is \emph{admissible} if its small values are
arbitrary and its big values are admissible. An element $f$ of a $\Pi$-code $p$ has \emph{bounded support} if for
every level $N\le\kappa(p)$ and every majorant $\chi$ there is an ordinal
$\beta^{(N)}_f(\chi)<\lambda_N$ such that for all $\alpha\ge\beta^{(N)}_f(\chi)$
and all arguments $G,G'$ with $\hat G,\hat G'\le\chi$: if
$r^N_\alpha G=r^N_\alpha G'$ then $r^N_\alpha f(G)=r^N_\alpha f(G')$; and,
recursively, the values $f(G)$ for $\hat G\le\chi$ have bounded support
with support functions bounded by a common function
$\beta^{(N)}_{f,\chi}(\cdot)<\lambda_N$, uniformly in $G$. Elements of
other codes have bounded support if their components do: the
components of a pair or a tree, the relation and base of a quotient
code, the families of a code, and some member of a class.
Three features are deliberate. Equality is demanded only \emph{after
retraction}: a function into codes, such as $x\mapsto\code\Sigma(X,x)$,
has no support in the literal sense, but is continuous in this one. The
condition is imposed at every level up to the value level $\kappa(p)$:
at $N=\kappa(p)$ it says that an element of $\Type_1\to\Type_0$ may
depend on its argument only through the argument's shadow below some
rank $<\lambda_0$; at $N<\kappa(p)$ it says that the level-$N$ shadow of
$f$'s output depends only on the level-$N$ shadow of its argument,
which prevents a non-term $F:\Type_N\to\Type_{N+1}$ from encoding
high-rank information about its argument in a low-rank part of its
output, where a level-$N$ function could read it. And the conditions at
$N>\kappa(p)$ are omitted because they follow from the one at
$N=\kappa(p)$ by compatibility of the retractions. Terms behave this
way at every level, since a body of value level $N$ can use a
higher-level argument only through propositions, through admissible
functions applied to it, or through a recursor along its actual
level-$\le N$ shape, as in Lemma~\ref{lem:stab}. The class $\mathcal A$ of \emph{admissible elements} is
the least class such that $f\in\mathcal A$ iff $f$ has bounded support
and $f\lhd M^{\mathcal A}_t(\rho,\tau)$, the $\mathcal A$-majorant of
Definition~\ref{def:relmaj}, for some term $t$ whose type has syntactic
value level at most $\kappa(p)$, as Definition~\ref{def:wtree} requires
of witnesses, some admissible
environment $\rho$ with values in decodings of codes in $\bigcup_iU_i$,
and $\tau$ the majorants witnessing the admissibility of $\rho$'s big
values. (Bounded support was added after the attempted proof of
Lemma~\ref{lem:transfer}; see Example~\ref{ex:probe}.) The \emph{model} decodes $c$ by
\[
\El(c):=\{f{\restriction}\El(c'):\ f\in\El^{\mathrm f}(c)\cap\mathcal A\}\quad\text{for }c=\code\Pi(c',G),
\]
and by $\El^{\mathrm f}(c)$ otherwise, with the components of pairs, trees,
and classes restricted likewise.
\end{definition}

\begin{remark}[free versus bound variables]\label{rem:freebound}
The choice of environments is embedded in the interpretation at two
places, which play different roles. A context variable $x{:}A$ ranges
over $\El(\sem A\rho)$, the model's decoding, so extending a context
adds an admissible value: the invariant of the soundness induction is
that free variables are admissible. A bound variable inside an
abstraction is never required to be admissible: the majorant clause
sums $M_b$ over all $a\in\El^{\mathrm f}(\sem A\rho)$, but only over those
with $\hat a\le\chi$, so the sum is controlled by the input majorant
rather than by admissibility. The full abstraction $\sem{\lambda x.b}\rho$
on all of $\El^{\mathrm f}(\sem A\rho)$ is used for one purpose only: its
restriction to $\El(\sem A\rho)$ is the model's element, and its
admissibility (Lemma~\ref{lem:termadm}) is what makes that restriction a
model element. Likewise the family $G$ of a code $\code\Sigma(c,G)$ is
the full family, and the rank of the code sums over all of $\El^{\mathrm f}(c)$;
Lemma~\ref{lem:tyform} bounds this sum through $M_{\lambda x.B}$ at
$\sigma'_c$, which majorizes every element of $\El^{\mathrm f}(c)$,
admissible or not. The level structure is what keeps the two roles from
colliding: if $A$ is small, $\El^{\mathrm f}(\sem A\rho)$ has only elements
of rank below $\lambda_0$, all admissible, so nothing non-admissible is
ever summed; if $A$ is big, the $\Sigma$ lives at the sort level of $A$
or above, and a sum over non-admissible elements, such as over a climbing
$F$ in $\lambda F.\,\Sigma n{:}\N.\,F^n(\N)$, reaches at most $\lambda_0$,
which is harmless below $\lambda_1$. The one place where such a sum could
hurt, a $\Sigma$ at $\Type_0$, is the place where the domain is small and
there is nothing to sum.
\end{remark}

The operator defining $\mathcal A$ is monotone: a larger $\mathcal A$
admits more environments and, by Definition~\ref{def:relmaj}, larger
majorants $M^{\mathcal A}_t$, and $\lhd$ is upward closed; so the least
fixed point exists. Among elements with bounded support it is downward closed under
majorization: if $g\lhd m\le M^{\mathcal A}_t(\rho,\tau)$ and $g$ has bounded support
then $g\in\mathcal A$; majorization alone does not suffice
(Example~\ref{ex:probe}). Every function on a small domain has bounded
support, since its arguments are fixed by $r_\alpha$ once $\alpha$
exceeds their ranks, and is admissible with the constant term
$\lambda x.\,p$ and a small parameter $p$ of sufficient rank as witness;
so small $\Pi$-codes decode fully. A small-valued function on a big
domain need not have bounded support.
Restricting to admissible arguments makes extensional equality of the
model coincide with set equality, since the model's quantifiers range
over $\El$.

\begin{lemma}[terms are admissible]\label{lem:termadm}
If $\Gamma\vdash t:T$ and $\rho$ is an admissible environment for
$\Gamma$ then $\sem t\rho\in\mathcal A$, with witness $(t,\rho,\tau)$.
\end{lemma}
\begin{proof}
$\sem t\rho\lhd M^{\mathcal A}_t(\rho,\tau)$ is the majorization clause
of $\mathcal A$; it holds by induction on $t$ because the model's
abstractions are restricted to admissible arguments, over which the
$\mathcal A$-majorant sums, and because the values of subterms lie in
the decodings of their types, by the induction on derivations of which
this lemma and Lemma~\ref{lem:tyform} are the two halves.
Bounded support is proved simultaneously with the commutation of
Lemma~\ref{lem:transfer}, by induction on $t$: the commutation
$r^N_\alpha\sem b\rho[x{:=}G]=\sem b^{\alpha}(r^N_\alpha\rho[x{:=}r^N_\alpha G])$,
valid for every level $N\le\kappa$ and every $\alpha$ above the leaves of $M_b(\rho[x{:=}G],\tau)$, the
reached ranks, and the supports of $\rho$'s elements at the reached
majorants, shows that $r^N_\alpha\sem{\lambda x.b}\rho(G)$ depends on $G$
only through $r^N_\alpha G$; and for $\hat G\le\chi$ these three bounds are
below $\lambda_N$, the first by closure, the second because reached
codes are majorized, the third by the hypothesis on $\rho$. So
$\beta(\chi)$ exists. The other constructs reduce to this one.
\end{proof}

So every rule that forms a \emph{term} is sound for free; soundness of
the rules that form \emph{types} is the one substantive lemma.

\begin{lemma}[type formation]\label{lem:tyform}
Let $\Gamma\vdash A:\Type_i$ and $\Gamma,x{:}A\vdash B:\Type_i$, and let
$\rho$ be admissible. Then $\sem{\Sigma x{:}A.B}\rho$, $\sem{\Pi x{:}A.B}\rho$,
$\sem{Wx{:}A.B}\rho\in U_i$, and likewise for quotients and for type
formers at $\Prop$.
\end{lemma}
\begin{proof}
By the induction on derivations, $c:=\sem A\rho\in U_i$, and, for every
type-valued proper subterm $t'$ of $B$ standing under binders
$\vec y{:}\vec C$ of $B$ (including $x$), the term
$\Sigma x{:}A.\,\Sigma\vec y{:}\vec C.\,t'$ is a smaller derivable type
at $\Type_i$; so its code lies in $U_i$ and the ranks of the values of
$t'$ under all admissible extensions of $\rho$ are bounded by an
ordinal $\delta_{t'}<\lambda_i$. Let $\delta<\lambda_i$ bound all of
these, together with $\rk(c)\cdot\omega$, the ranks of the level-$i$
values of $\rho$, and the level-$i$ supports $\beta^{(i)}_{\rho(F)}(\chi)$
of $\rho$'s $i$-big elements at the majorants $\chi$ reached, which are
below $\lambda_i$ by bounded support. Apply Lemma~\ref{lem:transfer} to
$(\lambda x.B,\rho,\tau)$ with $\alpha:=\delta+1$: with $\theta$ the
witness tree of $\rho$'s big variables and $\sigma$ the actual truth of
the model, $(r^i_\alpha\rho,\tau)$ is a witnessed
$(i,\alpha,\vec\lambda_{>i},\sigma)$-environment for $(\lambda x.B,\theta)$
and
\[
\ell_\alpha\bigl(M^{\mathcal A}_{\lambda x.B}(\rho,\tau)\bigr)\le
\ell\bigl(M_{\lambda x.B}(r^i_\alpha\rho,\tau)\bigr)\le
\nu^{i,0}_{\lambda x.B,\theta}(\alpha)\le\nu^i_{\lambda x.B,\theta}(\alpha)<\lambda_i,
\]
the second inequality because the higher extents $\vec\lambda_{>i}$ are
one choice of $\vec\alpha'$ in Definition~\ref{def:nu}, and the last by
closure. Now for $a\in\El(c)$, the family value
$G(a)=\sem B\rho[x{:=}a]$ is majorized by
$M^{\mathcal A}_B(\rho[x{:=}a],\tau)\le M^{\mathcal A}_{\lambda x.B}(\rho,\tau)(\hat a)$
by Lemma~\ref{lem:termadm}, so $\rk G(a)$ is below the ordinal part of
that value; and $\hat a$ is hereditarily below $\rk(c)\cdot\omega<\alpha$,
so that ordinal part is an $\alpha$-leaf. Hence
$\sup_{a\in\El(c)}\rk G(a)<\lambda_i$, uniformly, and the code
$\code Q(c,G{\restriction}\El(c))$, a pair of $c$ and a family whose
values on non-admissible arguments are the same terms' values and are
bounded in the same way, has rank below the limit $\lambda_i$.
Propositional formers produce truth values, which are in $U_0$.
\end{proof}

\begin{theorem}[conjectural]\label{thm:upper}
The structure of Definition~\ref{def:adm} is a model of $\CIC_\omega+\EM$,
so $\ZF\vdash\Con(\CIC_\omega+\EM)$; relativised to $\ZFh_{n+1}$ with
$U_i$ the given universes it gives $\ZFh_{n+1}\vdash\Con(\CIC_{n+2}+\EM)$.
\end{theorem}
\begin{proof}[Proof sketch]
Types are codes in the right universe by Lemma~\ref{lem:tyform}; terms
are admissible by Lemma~\ref{lem:termadm} and lie in the decodings of
their types because the full decoding is Werner's model of the term
rules; the computation rules, $\eta$, $\beta$, reflection, and proof
irrelevance hold because evaluation is set-theoretic and admissibility
is a property of the resulting set; cumulativity is literal; $\EM$ holds
since $\Prop$ has two elements; and $\bot$ decodes to $\emptyset$.
\end{proof}

\subsection{The transfer step}\label{sec:transfer}

Lemma~\ref{lem:tyform} needs to pass from an environment of the model,
whose universes have extent $\lambda_i$, to an environment of the
$\alpha$-model for some $\alpha<\lambda_i$. Values are not independent of
the extent, and a $\lambda$-model element of a big type, a function on
functions $U_0^{\lambda}\to U_0^{\lambda}$, is not an element of any
$\alpha$-decoding; so the passage must be the retraction of
Definition~\ref{def:retr}, and Example~\ref{ex:probe} shows that
majorization alone does not make it commute with evaluation.

\begin{lemma}[transfer]\label{lem:transfer}
Let $\rho$ be an admissible $\sigma$-environment of the model for
$\Gamma$, $t$ a term of value level at most $i$ with witness tree
$\theta$ for $\rho$'s big variables, and $\alpha<\lambda_i$ an ordinal
above the ranks of the level-$i$ values of $\rho$, above the ranks of
the values of all type-valued proper subterms of $t$ under all
admissible extensions of $\rho$, and above the level-$i$ supports
$\beta^{(i)}_{\rho(F)}(\chi)$ of the $i$-big elements of $\rho$ at all
majorants $\chi$ reached. Then $(r^i_\alpha\rho,\tau)$ is a witnessed
$(i,\alpha,\vec\lambda_{>i},\sigma)$-environment for $(t,\theta)$
(Definition~\ref{def:wenv}); for every proper subterm $t'$ of $t$ and
every admissible extension $\rho'$ of $\rho$ to the binders above $t'$,
$r^i_\alpha\sem{t'}^{\sigma}\rho'=\sem{t'}^{\sigma,(i,\alpha)}(r^i_\alpha\rho')$;
and
\[
\ell_\alpha\bigl(M^{\mathcal A}_t(\rho,\tau)\bigr)\le\ell\bigl(M_t(r^i_\alpha\rho,\tau)\bigr).
\]
The $i$-small elements of $\rho$ are fixed by $r^i_\alpha$ and need no
hypothesis.
\end{lemma}
The hypotheses are met below $\lambda_i$ by the argument in the proof of
Lemma~\ref{lem:tyform}. The conclusion is an inequality of
$\alpha$-leaves, not an equality of leaf suprema, for the reason given
after Definition~\ref{def:relmaj}; and the value commutation is claimed
for proper subterms only, since the root's own value may have rank
$\ge\alpha$, which is exactly what Lemma~\ref{lem:tyform} is about to
bound.
\begin{proof}[Proof sketch]
The commutation is proved by induction on $t'$, simultaneously for all
levels, for admissible $\rho'$; the ranks of all values met are below
$\alpha$ by hypothesis, so no universe clause of $r^i_\alpha$ collapses a
value to $\bot$.
\emph{Variables:} by definition of $r^i_\alpha\rho'$.
\emph{Application $f\,a$:} $r^i_\alpha\sem{f\,a}\rho'=r^i_\alpha(F(A))$ with
$F=\sem f\rho'$, $A=\sem a\rho'$; and $\sem{f\,a}^{\alpha}(r^i_\alpha\rho')=r^i_\alpha(F)(r^i_\alpha A)=r^i_\alpha(F(e^i_\alpha r^i_\alpha A))$
by induction. Now $r^i_\alpha A=r^i_\alpha(e^i_\alpha r^i_\alpha A)$ since $r^i_\alpha e^i_\alpha=\mathrm{id}$,
and $\hat A$ is a reached majorant with $\alpha\ge\beta_F(\hat A)$, so
$r^i_\alpha F(A)=r^i_\alpha F(e^i_\alpha r^i_\alpha A)$ by bounded support of $F$,
which holds because $F$ is either an element of $\rho'$ or, by the
simultaneous induction of Lemma~\ref{lem:termadm}, the value of a term in
an admissible environment.
\emph{Abstraction $\lambda x.b$:} the model's abstraction is the
restriction to $\El(\sem A\rho')$ of $a\mapsto\sem b\rho'[x{:=}a]$, and
$r^i_\alpha$ of it is $a'\mapsto r^i_\alpha\sem b\rho'[x{:=}e^i_\alpha a']$,
which by induction is $a'\mapsto\sem b^{\alpha}(r^i_\alpha\rho'[x{:=}a'])$
on the retracted admissible arguments, using $r^i_\alpha e^i_\alpha=\mathrm{id}$.
\emph{Data of small type and their eliminators} ($\mathbf n$, $\N$,
$\Sigma$ and $=$ at small types): data of rank below $\alpha$ is fixed
by $r^i_\alpha$, so the branch selected and the recursion performed are
the same on both sides, and the step functions commute by the
application case.
\emph{$W$-recursion:} let $V(t):=\sem{\mathsf{ind}^W_Cf}\rho'(t)=F(a)(d)(\mathrm{rec}_t)$
for $t=\mathsf{tree}\,a\,d$, with $F:=\sem f\rho'$ and
$\mathrm{rec}_t(b):=V(d(b))$, and let $V^\alpha$ be the same in the cut
structure. By induction on the tree,
$r^i_\alpha V(t)=(r^i_\alpha F)(r^i_\alpha a)(r^i_\alpha d)(r^i_\alpha\mathrm{rec}_t)$
by three uses of the application case, which needs $\alpha$ above the
supports of $F$, $F(a)$ and $F(a)(d)$ at the majorants $\hat a$,
$\hat d$, $\widehat{\mathrm{rec}_t}$, all reached; and
$(r^i_\alpha\mathrm{rec}_t)(b')=r^i_\alpha V(d(e^i_\alpha b'))=V^\alpha(r^i_\alpha d(e^i_\alpha b'))=V^\alpha((r^i_\alpha d)(b'))$
by the induction hypothesis at the subtree $d(e^i_\alpha b')$, so
$r^i_\alpha\mathrm{rec}_t=\mathrm{rec}^\alpha_{r^i_\alpha t}$ and
$r^i_\alpha V(t)=V^\alpha(r^i_\alpha t)$. The retraction of $t$ is well
defined because the branching family of the $W$-code is the value of
$\lambda x.B$ in an admissible environment and has bounded support.
\emph{$\mathsf{Acc}$-recursion:} let $V(x):=\sem{\mathsf{ind}^{\mathsf{Acc}}_Cf}\rho'(x)(\star)=F(x)(\mathrm{rec}_x)$
for accessible $x$, with $\mathrm{rec}_x(y)(\star):=V(y)$ for
$\sem R\rho'(y,x)=\top$. The relation of the cut structure is
$\sem R^\alpha(r^i_\alpha\rho')(y',x')=\sem R\rho'(e^i_\alpha y',e^i_\alpha x')$
by the commutation for the subterm $R$, and
$\sem R\rho'(e^i_\alpha y',e^i_\alpha r^i_\alpha x)=\sem R\rho'(e^i_\alpha y',x)$
by bounded support of $\sem R\rho'$ (a term value) at $\alpha$ above its
support; so the predecessors of $r^i_\alpha x$ in the cut structure are
the $y'$ with $e^i_\alpha y'$ a predecessor of $x$, and $r^i_\alpha x$ is
accessible when $x$ is, by induction along the accessible part. Then,
by the same induction,
$r^i_\alpha V(x)=(r^i_\alpha F)(r^i_\alpha x)(r^i_\alpha\mathrm{rec}_x)$
by the application case, and
$(r^i_\alpha\mathrm{rec}_x)(y')(\star)=r^i_\alpha V(e^i_\alpha y')=V^\alpha(y')$
by the induction hypothesis at the predecessor $e^i_\alpha y'$, so
$r^i_\alpha\mathrm{rec}_x=\mathrm{rec}^\alpha_{r^i_\alpha x}$ and
$r^i_\alpha V(x)=V^\alpha(r^i_\alpha x)$. The truth values of $R$ are read
under the same $\sigma$ on both sides.
\emph{Quotients:} $\sem{\mathsf{ind}^/_Cf\,h}\rho'([x])=F(x)$ for any
member $x$; $r^i_\alpha F(x)=(r^i_\alpha F)(r^i_\alpha x)$ by the application
case, and $r^i_\alpha x$ is a member of $r^i_\alpha[x]=[r^i_\alpha x]_{R'}$,
so this is the cut structure's value at $r^i_\alpha[x]$. The class map is
well defined by the remark after Definition~\ref{def:retr}, the
relation being a term value with bounded support.
\emph{Type formers:} $r^i_\alpha$ on codes is defined to commute with the
formers, and the universe clause of $r^i_\alpha$ is the identity on the
reached level-$i$ codes because their ranks are below $\alpha$; higher
universes are the true ones in the target structure, so no collapsing
occurs there.
\emph{The inequality of $\alpha$-leaves}, by induction on $t$ over the
clauses of Definition~\ref{def:hyb}: for the abstraction clause at an
input $\chi$ hereditarily below $\alpha$, the $\mathcal A$-sum ranges
over admissible $a$ with $\hat a\le\chi$; each such $a$ maps to
$r^i_\alpha a$, an element of the cut decoding with $\widehat{r^i_\alpha a}\le\hat a\le\chi$,
and $\ell_\alpha(M^{\mathcal A}_b(\rho'[x{:=}a],\tau))\le\ell(M_b(r^i_\alpha\rho'[x{:=}r^i_\alpha a],\tau))$
by induction, since $\rho'[x{:=}a]$ is admissible; so the
$\alpha$-leaves of the $\mathcal A$-sum are bounded by leaves of the
full sum in the cut structure. Inputs not hereditarily below $\alpha$
contribute no $\alpha$-leaves. In the other clauses the majorant is
assembled from majorants of subterms at canonical majorants of reached
values, and reached values commute with $r^i_\alpha$ and have equal
canonical majorants on both sides (their ranks being below $\alpha$),
so the clauses are compared termwise.
Finally $(r^i_\alpha\rho,\tau)$ is a witnessed environment because
$r^i_\alpha\rho$ is a $\sigma$-environment (types commute with $r^i_\alpha$
by the type-former case) and each $r^i_\alpha\rho(F)$ is majorized by
$\tau(F)=M^{\mathcal A}_{t_F}(\rho_F,\tau_F)$, whose $\alpha$-leaves are
bounded by those of $M_{t_F}(r^i_\alpha\rho_F,\tau_F)$ by the same
argument applied to the witness trees, with
$(r^i_\alpha\rho_F,\tau_F)$ witnessed by the induction on $\theta$; here
majorization of $r^i_\alpha\rho(F)$ by the cut majorant needs only its
values at inputs hereditarily below $\alpha$, since $r^i_\alpha\rho(F)$
takes arguments of rank below $\alpha$.
\end{proof}

The proof is a logical relation between the two structures, with
bounded support as the single semantic hypothesis on the environment;
the $W$, $\mathsf{Acc}$, and quotient cases are now written, and each
reduces to the application case together with the well-definedness of
the retraction on the data, which is bounded support of the families
and relations involved. The
last step, that $\tau(F)$ may be replaced in the witnessed environment
by a majorant agreeing with it below $\alpha$, needs Definition~\ref{def:wenv}
to be read with $\tau(F)$ replaced by $M_{t_F}(r^i_\alpha\rho_F,\tau_F)$,
which is what the definition says when the recursion is unwound in the
cut structure.

\begin{remark}[a gap in the abstraction case]\label{rem:gap}
The abstraction case of the inequality pairs each admissible argument
$a$, $\hat a\le\chi$, with its retraction $r^i_\alpha a$ and appeals to
the induction hypothesis for the environment $\rho'[x{:=}a]$. But the
hypotheses of the lemma include that $\alpha$ lies above the supports
of the environment's elements, and the supports of the admissible $a$
with $\hat a\le\chi$ are not bounded below $\lambda_i$, because the
canonical majorant does not see the support. Let
$a_\gamma:(\Type_i\to\Type_i)\to\N$ be $0$ at $G$ if $G$ is the identity
on codes of rank below $\gamma$ and $1$ otherwise. Each $a_\gamma$ is
admissible, with support $\gamma$ and majorant $2$; and for
$\alpha<\gamma$, $(r^i_\alpha a_\gamma)(r^i_\alpha\mathrm{id})=a_\gamma(e^i_\alpha r^i_\alpha\mathrm{id})=1\ne0=a_\gamma(\mathrm{id})$.
So a body that branches on $x\,(\lambda X.X)$ takes different branches
at $\rho'[x{:=}a_\gamma]$ and at its retraction, the commutation fails
for this summand, and the summand's majorant is compared with the wrong
branch. In this example the inequality survives, because the cut sum
also contains the constant functions, which select either branch; in
general one needs, for each summand of the model's sum, some element of
the cut structure that reproduces the summand's values at all the
arguments the body applies it to, and when two such arguments have the
same retraction and different values (the identity, and a bound
variable equal to the identity below $\alpha$) the retraction of the
second must be replaced too, which changes the other leaves that depend
on it. We have not found a counterexample to the inequality, and the
majorizability heuristic says it should hold, since no leaf of a term
majorant depends on the support of an input; but the proof above does
not establish it.

The repair we expect to work is to make the support a component of the
majorant: a function part carries, besides its values, an ordinal
modulus, $f\lhd m$ requires $\beta^{(N)}_f(\chi)\le m_{\mathrm s}(\chi)$,
the hybrid majorant of an abstraction computes the modulus from the
moduli of the environment and the leaves reached, exactly as in the
proof of Lemma~\ref{lem:termadm}, and the leaf supremum, hence the
closure functions, include the moduli. Then the sum over
$\hat a\le\chi$ ranges only over arguments of support at most
$\chi_{\mathrm s}$, which is below $\alpha$ for the inputs that matter,
every summand commutes with the retraction, and the elements of a small
code, on which $r^i_\alpha$ is the identity, have modulus below
$\alpha$ as Lemma~\ref{lem:tyform} needs. This is how majorants of
continuous functionals carry moduli of continuity. It moves the
retraction between cut structures into the first phase, which is
harmless since it never used the true extents, and it replaces the
separate bounded-support clause of Definition~\ref{def:adm}. It has not
been carried out.
\end{remark}

\begin{lemma}[bounded support is preserved]\label{lem:bspres}
In an admissible environment $\rho'$, the values of the eliminators
$\mathsf{ind}^W_Cf$, $\mathsf{ind}^{\mathsf{Acc}}_Cf$, and $\mathsf{ind}^/_Cf\,h$
have bounded support at every level $N\le\kappa$ of the motive.
\end{lemma}
\begin{proof}
Let $V$ be the value, a function on trees, accessible elements, or
classes. Fix $N$ and an input majorant $\chi$. The commutation cases of
Lemma~\ref{lem:transfer} show $r^N_\alpha V(t)=V^\alpha(r^N_\alpha t)$ for
every $\alpha$ above the supports of $F=\sem f\rho'$ and of its
iterated values at the majorants reached along $t$, and above the
support of the branching family or relation. For $t$ with
$\hat t\le\chi$ the reached majorants are bounded in terms of $\chi$:
the labels' majorants are components of $\hat t\le\chi$, the subtree
functions' likewise, and the recursive functions' majorants are bounded
by $M_{\mathsf{ind}^W_Cf}(\rho',\tau)(\chi)$, a fixed majorant; for
$\mathsf{Acc}$ the predecessors' majorants are below $\chi$ (they are
elements of the same code) and the recursive functions' by
$M_{\mathsf{ind}^{\mathsf{Acc}}_Cf}(\rho',\tau)(\chi)$; for quotients the
members' majorants are below $\chi$ by the definition of $\hat q$. Since
the support functions of $F$ and of its values are monotone (the
condition at $\chi$ subsumes the condition at every $\chi'\le\chi$)
and uniform in the curried arguments by Definition~\ref{def:adm},
their values at these finitely many bounds give one ordinal
$\beta(\chi)<\lambda_N$, and for $\alpha\ge\beta(\chi)$ the equation
$r^N_\alpha V(t)=V^\alpha(r^N_\alpha t)$ shows that $r^N_\alpha V(t)$
depends on $t$ only through $r^N_\alpha t$. The values $V(t)$ have
bounded support with bounds uniform in $t$ for $\hat t\le\chi$ by the
same argument one argument deeper, using that they are values of $F$
at inputs bounded in terms of $\chi$.
\end{proof}

\paragraph{Remaining checks.}
The degree bookkeeping for recursors on actual trees is uniform: all
values lie in decodings of subcodes of the motive, so they have degree
$\le\rk c+1$ for the code $c$ of the result type. Restriction to
admissible arguments is compatible with abstraction, since the model's
element $\sem{\lambda x.b}\rho{\restriction}\El(\sem A\rho)$ is by
Lemma~\ref{lem:termadm} the restriction of an admissible full function,
and with $\eta$ and $\beta$, by restricting both sides. Reflection holds
because $\sem a\rho=\sem{a'}\rho$ as sets gives equal canonical
majorants. For the relativisation to $\ZFh_{n+1}$, the closure functions
$\nu_{t,\theta}$, the hybrid majorants, the truth assignments, and the
retraction are all defined by recursion on syntax over $\omega$ and by
set recursion, hence are $\Lang$-terms, so the given universes $U_i$ are
closed under them; the same open lemma is all that is needed there.

\subsection{The two examples that shaped the design}\label{sec:examples}

\begin{example}[heterogeneous fibres]\label{ex:hetero}
Let $B(0):=\N$, $B(1):=\Type_2$ and $f:=\mathsf{ind}^{\mathbf 2}\,5\,c$
with $c\in U_2$ of rank $\ge\lambda_0$. A purely syntactic majorant of
$f$ at an argument majorant $2$, which covers both $0$ and $1$, is the
join of both branches, so $\lambda Y{:}\Type_0.\,T(f\,(\mathsf{swap}\,b'))$
with $\rho(b')=1$, whose values are all small, would receive a majorant
above $\lambda_0$; and a pointwise structured condition
``$m(\chi)$ admissible for every fibre $G(a)$ with $a\lhd\chi$'' is
unsatisfiable at $\chi=2$. Branching on the actual value of a data-typed
argument removes the join; and a big branch can reach a small position
only in a context that is empty in the model, such as $e{:}b=0$ at $b=1$,
where the abstraction supremum over the actual elements of $\El(\bot)$
is empty.
\end{example}

\begin{example}[a climbing function]\label{ex:climb}
$\lambda_0$ has cofinality $\omega$; fix a cofinal sequence $(\beta_n)$
and let $\mathrm{next}(\beta)$ be the least $\beta_n>\beta$, and
$F(X):=$ a code of rank $\mathrm{next}(\rk X)$. Any \emph{pointwise}
admissibility condition, such as ``$\hat F(\chi)<\lambda_0$ for
$\chi<\lambda_0$ and $\hat F(\lambda_0)<\lambda_1$,'' accepts $F$, yet
$\Sigma n{:}\N.\,F^n(\N):\Type_0$ denotes a code of rank $\lambda_0$. No
condition on the values of $F$ can stop an $\omega$-iteration from
climbing an ordinal of cofinality $\omega$; what stops it in
Section~\ref{sec:rank} is that functions there are terms and $\lambda_0$
is closed under the iterate of every term majorant. Definition~\ref{def:adm}
imports exactly this: $F$ is not admissible because no term majorant
dominates $\mathrm{next}$, since otherwise the $\omega$-iterate of that
majorant, which is the majorant of an $\mathsf{ind}^\N$ term, would
reach $\lambda_0$, contradicting closure.
\end{example}

\paragraph{Status.}
Verified at the level of detail above: the codes and full decoding are
independent of the $\lambda_i$ except through the universe constants,
whose rank must be $\lambda_i+1$ for the family domains to cover
$U_i$ (an eighth correction, found when ordering the definitions);
hybrid majorants are sound and
model-independent; the closure ordinals, now with proofs
(Lemmas~\ref{lem:sigmaloc}, \ref{lem:envset}, \ref{lem:nuset},
\ref{lem:closedex}), after a ninth correction: the closure functions
were not monotone in $\alpha$ and not continuous, since big values on
$U_N^\beta$ are not elements at $\alpha>\beta$, so monotonicity is now
built into the definition and continuity, which was never used, is
dropped;
the monotonicity of the admissibility operator; Lemmas~\ref{lem:termadm}
and \ref{lem:tyform}; Examples~\ref{ex:set}, \ref{ex:hetero}, \ref{ex:climb}.
The level recursion is bottom-up only because the auxiliary models take
arbitrary higher extents and Lemma~\ref{lem:stab} bounds the resulting
supremum; without that, stage $N$ would depend on $\lambda_{N+1}$ through
the domains of $N$-big variables of higher level, and a witness-majorized
$Y{:}\Type_{N+1}$ of large rank would be collapsed by the retraction and
break the commutation on $Y\to Y$. Lemma~\ref{lem:stab} is now proved,
after being restricted to roots of value level $\le N$ (without which
it is false for the term $\Type_{N+1}$) and after Lemma~\ref{lem:auxrank},
a rank lemma for the auxiliary structure, was separated from it: the
values of $W$- and $\mathsf{Acc}$-recursions are not majorized in the
auxiliary structure and are bounded instead by the height of the part
of the datum that a term can inspect. With that, and with the
eliminator clauses of Definition~\ref{def:hyb} now written out, the
first phase has no conjecture and no informal definition left.
Lemma~\ref{lem:transfer} now has a proof sketch, but only after a
seventh correction: majorization alone does not make the retraction a
homomorphism (Example~\ref{ex:probe}), so admissibility now also requires
bounded support, formulated as continuity with respect to the
retraction and bounded at the value level, and Lemma~\ref{lem:termadm}
is proved by a simultaneous induction with the commutation. Pushing the
condition through \S\ref{sec:adm} changed two statements: downward
closure under majorization holds only among elements with bounded
support, and small-valued functions on big domains are no longer
automatically admissible, while functions on small domains are.
A tenth correction changed the transfer's conclusion: ``the same leaf
supremum'' is false already for $\lambda X{:}\Type_i.X$ (the remark
after Definition~\ref{def:relmaj}), and the model's majorants must sum
over admissible arguments only, so the lemma now compares the
$\alpha$-leaves of the model's $\mathcal A$-majorant with the leaves of
the full majorant in the cut structure, and Lemma~\ref{lem:tyform}
obtains its cut $\alpha$ uniformly in the family argument from the
induction hypothesis applied to the smaller types
$\Sigma x{:}A.\Sigma\vec y.\,t'$ formed from the proper subterms of $B$,
which removes a circularity in the earlier proof (the cut had to lie
above the ranks it was meant to bound). The $W$, $\mathsf{Acc}$, and
quotient cases of the commutation are now written, and
Lemma~\ref{lem:bspres} preserves bounded support through those
eliminators. Writing them exposed that the retraction of
Definition~\ref{def:retr} is well defined on pairs, trees, and classes
only when the families and relations involved have bounded support
(the fibre at $r a$ must not depend on the choice of $a$, and
$e$ on classes needs $x\sim_R e\,r\,x$), so the definition is now
stated for admissible codes, the quotient relation is pushed forward
rather than pulled back, and bounded support is required uniformly in
curried arguments; the apparent circularity between the retraction and
bounded support is not one, the retraction being a set map defined
without conditions. We do not claim the theorem. The value commutation
of Lemma~\ref{lem:transfer} is proved case by case, but its inequality
of $\alpha$-leaves has a gap in the abstraction case
(Remark~\ref{rem:gap}): the model's majorant sums over admissible
arguments whose supports are unbounded below $\lambda_i$, since the
canonical majorant does not see the support, and such an argument does
not commute with a retraction at a fixed cut. The expected repair,
making the support modulus a component of the majorant, is described
there and not carried out. Checking that repair against the examples
produced an eleventh correction, which supersedes it: the closure
functions, being computed in a cut structure, saturate at the cut, so
closure under them is nearly vacuous, $\lambda_N$ is the least ordinal
closed under Hartogs numbers, and Lemma~\ref{lem:tyform} is false for
the $\omega$-iterate of a Hartogs-raising type operator
(Example~\ref{ex:sat}). This is a defect of the definition of
$\lambda_N$ and not of any proof. The repair of Remark~\ref{rem:nosat},
a supremum over all cuts of the leaves at small inputs, would restore
the argument of Example~\ref{ex:climb}, which as it stands is unfounded,
and would make the model's own structure one of those supped over, so
that type formation follows from soundness of the majorants and closure
with no transfer, retraction, or support condition. What it needs is a
cut-stability lemma, which is open. The present state is therefore:
Sections~\ref{sec:codes} and \ref{sec:hybrid} stand; the lemmas of
\S\ref{sec:closure} are proved but about the wrong ordinals;
\S\ref{sec:adm}--\ref{sec:transfer} describe a structure that is not a
model and machinery that the repair would remove.

{\small
\bibliographystyle{plain}
\bibliography{references}}

\end{document}
